%% main.tex -- CS 259 term paper (cs259.cls reproduces the course template)
%% Compile with pdfLaTeX + BibTeX (Overleaf default); see README_OVERLEAF.md.
\documentclass{cs259}
%% \documentclass[siunitx]{cs259}   % uncomment to use \SI{}{} and \num{}

\graphicspath{{figures/}}
\hypersetup{
  pdftitle  = {Performance Evaluation of Hybrid CPU-GPU Parallelism for Deep Learning Using OpenMP and CUDA},
  pdfauthor = {Sunnihith Reddy Induru}}

\title{Performance Evaluation of Hybrid CPU-GPU Parallelism\\
       for Deep Learning Using OpenMP and CUDA}
\author{Sunnihith Reddy Induru}
\affiliation{Computer Science Department\\
             San Jose State University\\
             San Jose, CA 95192}
\email{sunnihithreddy.induru@sjsu.edu}

\newcommand{\cudabase}{CUDA (1-thread preprocessing)}
\begin{document}
%% cs259ieee.bst does not sort. The course template wants a numbered list ordered alphabetically by first author,
%% so every reference is cited once here, in alphabetical order, before the first \cite in the text.
%% (When adding a reference, insert its key at the right place in this list.)
\makeatletter
\def\bstctlcite#1{\@bsphack\@for\@citeb:=#1\do{%
  \edef\@citeb{\expandafter\@firstofone\@citeb}\if@filesw\immediate\write\@auxout{\string\citation{\@citeb}}\fi}\@esphack}
\makeatother
\bstctlcite{IEEEexample:BSTcontrol}% BibTeX style settings (refs.bib)
\nocite{%
  amdahl1967validity,dagum1998openmp,gregg2011data,gustafson1988reevaluating,%
  harris2012metrics,harris2012transfers,harris2012overlap,he2015delving,hill2008amdahl,%
  hoefler2015scientific,induru2026repository,jouppi2017tpu,karp1990measuring,kingma2015adam,%
  lecun1998gradient,lee2010debunking,mittal2015survey,nickolls2008scalable,nvidia2026cudaguide,%
  openmparb2021specification,sze2017efficient,williams2009roofline,yoo2003slurm}
% Generated by analysis/generate_numbers.py from results/raw -- do not edit by hand.
\providecommand{\val}[1]{\csname v:#1\endcsname}
\expandafter\gdef\csname v:Mac:accuracy\endcsname{97.97}
\expandafter\gdef\csname v:Mac:amdahl:ceil:1\endcsname{1.0}
\expandafter\gdef\csname v:Mac:amdahl:ceil:1024\endcsname{11.2}
\expandafter\gdef\csname v:Mac:amdahl:ceil:128\endcsname{6.3}
\expandafter\gdef\csname v:Mac:amdahl:ceil:256\endcsname{7.9}
\expandafter\gdef\csname v:Mac:amdahl:ceil:32\endcsname{3.1}
\expandafter\gdef\csname v:Mac:amdahl:ceil:4096\endcsname{12.0}
\expandafter\gdef\csname v:Mac:amdahl:ceil:64\endcsname{4.4}
\expandafter\gdef\csname v:Mac:amdahl:ceil:8\endcsname{1.3}
\expandafter\gdef\csname v:Mac:amdahl:f:1\endcsname{1.000}
\expandafter\gdef\csname v:Mac:amdahl:f:1024\endcsname{0.089}
\expandafter\gdef\csname v:Mac:amdahl:f:128\endcsname{0.160}
\expandafter\gdef\csname v:Mac:amdahl:f:256\endcsname{0.126}
\expandafter\gdef\csname v:Mac:amdahl:f:32\endcsname{0.319}
\expandafter\gdef\csname v:Mac:amdahl:f:4096\endcsname{0.083}
\expandafter\gdef\csname v:Mac:amdahl:f:64\endcsname{0.228}
\expandafter\gdef\csname v:Mac:amdahl:f:8\endcsname{0.797}
\expandafter\gdef\csname v:Mac:amdahl:r2:1\endcsname{-3.51}
\expandafter\gdef\csname v:Mac:amdahl:r2:1024\endcsname{0.96}
\expandafter\gdef\csname v:Mac:amdahl:r2:128\endcsname{0.88}
\expandafter\gdef\csname v:Mac:amdahl:r2:256\endcsname{0.92}
\expandafter\gdef\csname v:Mac:amdahl:r2:32\endcsname{0.51}
\expandafter\gdef\csname v:Mac:amdahl:r2:4096\endcsname{0.97}
\expandafter\gdef\csname v:Mac:amdahl:r2:64\endcsname{0.73}
\expandafter\gdef\csname v:Mac:amdahl:r2:8\endcsname{-0.08}
\expandafter\gdef\csname v:Mac:compiler\endcsname{Apple LLVM 21.0.0 (clang-2100.3.34.2)}
\expandafter\gdef\csname v:Mac:cpu\endcsname{Apple M4}
\expandafter\gdef\csname v:Mac:cpuThreads\endcsname{10}
\expandafter\gdef\csname v:Mac:cpubest:e2e:1\endcsname{135.2}
\expandafter\gdef\csname v:Mac:cpubest:e2e:1024\endcsname{25.76}
\expandafter\gdef\csname v:Mac:cpubest:e2e:128\endcsname{36.17}
\expandafter\gdef\csname v:Mac:cpubest:e2e:256\endcsname{31.75}
\expandafter\gdef\csname v:Mac:cpubest:e2e:32\endcsname{47.24}
\expandafter\gdef\csname v:Mac:cpubest:e2e:4096\endcsname{24.92}
\expandafter\gdef\csname v:Mac:cpubest:e2e:64\endcsname{42.01}
\expandafter\gdef\csname v:Mac:cpubest:e2e:8\endcsname{80.51}
\expandafter\gdef\csname v:Mac:cpubest:impl:1\endcsname{OpenMP}
\expandafter\gdef\csname v:Mac:cpubest:impl:1024\endcsname{OpenMP}
\expandafter\gdef\csname v:Mac:cpubest:impl:128\endcsname{OpenMP}
\expandafter\gdef\csname v:Mac:cpubest:impl:256\endcsname{OpenMP}
\expandafter\gdef\csname v:Mac:cpubest:impl:32\endcsname{OpenMP}
\expandafter\gdef\csname v:Mac:cpubest:impl:4096\endcsname{OpenMP}
\expandafter\gdef\csname v:Mac:cpubest:impl:64\endcsname{OpenMP}
\expandafter\gdef\csname v:Mac:cpubest:impl:8\endcsname{OpenMP}
\expandafter\gdef\csname v:Mac:cudaRuntime\endcsname{none}
\expandafter\gdef\csname v:Mac:cv:max\endcsname{15.0}
\expandafter\gdef\csname v:Mac:cv:median\endcsname{0.6}
\expandafter\gdef\csname v:Mac:gpu\endcsname{none}
\expandafter\gdef\csname v:Mac:micro:cpu_triad_bandwidth:1\endcsname{90.03}
\expandafter\gdef\csname v:Mac:micro:cpu_triad_bandwidth:10\endcsname{86.46}
\expandafter\gdef\csname v:Mac:micro:cpu_triad_bandwidth:2\endcsname{94.23}
\expandafter\gdef\csname v:Mac:micro:cpu_triad_bandwidth:4\endcsname{92.07}
\expandafter\gdef\csname v:Mac:micro:cpu_triad_bandwidth:6\endcsname{93.44}
\expandafter\gdef\csname v:Mac:micro:cpu_triad_bandwidth:8\endcsname{92.77}
\expandafter\gdef\csname v:Mac:micro:openmp_fork_join:1\endcsname{0.05}
\expandafter\gdef\csname v:Mac:micro:openmp_fork_join:10\endcsname{21.98}
\expandafter\gdef\csname v:Mac:micro:openmp_fork_join:2\endcsname{3.89}
\expandafter\gdef\csname v:Mac:micro:openmp_fork_join:4\endcsname{7.25}
\expandafter\gdef\csname v:Mac:micro:openmp_fork_join:6\endcsname{13.05}
\expandafter\gdef\csname v:Mac:micro:openmp_fork_join:8\endcsname{19.94}
\expandafter\gdef\csname v:Mac:omp:best:e2e:1\endcsname{135.2}
\expandafter\gdef\csname v:Mac:omp:best:e2e:1024\endcsname{25.76}
\expandafter\gdef\csname v:Mac:omp:best:e2e:128\endcsname{36.17}
\expandafter\gdef\csname v:Mac:omp:best:e2e:256\endcsname{31.75}
\expandafter\gdef\csname v:Mac:omp:best:e2e:32\endcsname{47.24}
\expandafter\gdef\csname v:Mac:omp:best:e2e:4096\endcsname{24.92}
\expandafter\gdef\csname v:Mac:omp:best:e2e:64\endcsname{42.01}
\expandafter\gdef\csname v:Mac:omp:best:e2e:8\endcsname{80.51}
\expandafter\gdef\csname v:Mac:omp:best:eff:1\endcsname{1.04}
\expandafter\gdef\csname v:Mac:omp:best:eff:1024\endcsname{0.54}
\expandafter\gdef\csname v:Mac:omp:best:eff:128\endcsname{0.38}
\expandafter\gdef\csname v:Mac:omp:best:eff:256\endcsname{0.44}
\expandafter\gdef\csname v:Mac:omp:best:eff:32\endcsname{0.74}
\expandafter\gdef\csname v:Mac:omp:best:eff:4096\endcsname{0.56}
\expandafter\gdef\csname v:Mac:omp:best:eff:64\endcsname{0.83}
\expandafter\gdef\csname v:Mac:omp:best:eff:8\endcsname{0.43}
\expandafter\gdef\csname v:Mac:omp:best:speedup:1\endcsname{1.04}
\expandafter\gdef\csname v:Mac:omp:best:speedup:1024\endcsname{5.40}
\expandafter\gdef\csname v:Mac:omp:best:speedup:128\endcsname{3.83}
\expandafter\gdef\csname v:Mac:omp:best:speedup:256\endcsname{4.37}
\expandafter\gdef\csname v:Mac:omp:best:speedup:32\endcsname{2.95}
\expandafter\gdef\csname v:Mac:omp:best:speedup:4096\endcsname{5.56}
\expandafter\gdef\csname v:Mac:omp:best:speedup:64\endcsname{3.30}
\expandafter\gdef\csname v:Mac:omp:best:speedup:8\endcsname{1.72}
\expandafter\gdef\csname v:Mac:omp:best:thr:1\endcsname{1}
\expandafter\gdef\csname v:Mac:omp:best:thr:1024\endcsname{10}
\expandafter\gdef\csname v:Mac:omp:best:thr:128\endcsname{10}
\expandafter\gdef\csname v:Mac:omp:best:thr:256\endcsname{10}
\expandafter\gdef\csname v:Mac:omp:best:thr:32\endcsname{4}
\expandafter\gdef\csname v:Mac:omp:best:thr:4096\endcsname{10}
\expandafter\gdef\csname v:Mac:omp:best:thr:64\endcsname{4}
\expandafter\gdef\csname v:Mac:omp:best:thr:8\endcsname{4}
\expandafter\gdef\csname v:Mac:omp:best:tput:1\endcsname{73{,}943}
\expandafter\gdef\csname v:Mac:omp:best:tput:1024\endcsname{388{,}159}
\expandafter\gdef\csname v:Mac:omp:best:tput:128\endcsname{276{,}478}
\expandafter\gdef\csname v:Mac:omp:best:tput:256\endcsname{314{,}995}
\expandafter\gdef\csname v:Mac:omp:best:tput:32\endcsname{211{,}675}
\expandafter\gdef\csname v:Mac:omp:best:tput:4096\endcsname{401{,}235}
\expandafter\gdef\csname v:Mac:omp:best:tput:64\endcsname{238{,}045}
\expandafter\gdef\csname v:Mac:omp:best:tput:8\endcsname{124{,}202}
\expandafter\gdef\csname v:Mac:omp:e2e:1024:1\endcsname{139.2}
\expandafter\gdef\csname v:Mac:omp:e2e:1024:10\endcsname{25.76}
\expandafter\gdef\csname v:Mac:omp:e2e:1024:2\endcsname{71.22}
\expandafter\gdef\csname v:Mac:omp:e2e:1024:4\endcsname{36.26}
\expandafter\gdef\csname v:Mac:omp:e2e:1024:6\endcsname{35.54}
\expandafter\gdef\csname v:Mac:omp:e2e:1024:8\endcsname{27.62}
\expandafter\gdef\csname v:Mac:omp:e2e:128:1\endcsname{139.0}
\expandafter\gdef\csname v:Mac:omp:e2e:128:10\endcsname{36.17}
\expandafter\gdef\csname v:Mac:omp:e2e:128:2\endcsname{72.65}
\expandafter\gdef\csname v:Mac:omp:e2e:128:4\endcsname{39.07}
\expandafter\gdef\csname v:Mac:omp:e2e:128:6\endcsname{43.58}
\expandafter\gdef\csname v:Mac:omp:e2e:128:8\endcsname{36.26}
\expandafter\gdef\csname v:Mac:omp:e2e:1:1\endcsname{135.2}
\expandafter\gdef\csname v:Mac:omp:e2e:1:10\endcsname{1{,}092.2}
\expandafter\gdef\csname v:Mac:omp:e2e:1:2\endcsname{265.3}
\expandafter\gdef\csname v:Mac:omp:e2e:1:4\endcsname{352.8}
\expandafter\gdef\csname v:Mac:omp:e2e:1:6\endcsname{623.9}
\expandafter\gdef\csname v:Mac:omp:e2e:1:8\endcsname{863.1}
\expandafter\gdef\csname v:Mac:omp:e2e:256:1\endcsname{139.1}
\expandafter\gdef\csname v:Mac:omp:e2e:256:10\endcsname{31.75}
\expandafter\gdef\csname v:Mac:omp:e2e:256:2\endcsname{71.56}
\expandafter\gdef\csname v:Mac:omp:e2e:256:4\endcsname{37.47}
\expandafter\gdef\csname v:Mac:omp:e2e:256:6\endcsname{37.99}
\expandafter\gdef\csname v:Mac:omp:e2e:256:8\endcsname{32.12}
\expandafter\gdef\csname v:Mac:omp:e2e:32:1\endcsname{139.0}
\expandafter\gdef\csname v:Mac:omp:e2e:32:10\endcsname{63.86}
\expandafter\gdef\csname v:Mac:omp:e2e:32:2\endcsname{79.55}
\expandafter\gdef\csname v:Mac:omp:e2e:32:4\endcsname{47.24}
\expandafter\gdef\csname v:Mac:omp:e2e:32:6\endcsname{63.75}
\expandafter\gdef\csname v:Mac:omp:e2e:32:8\endcsname{57.29}
\expandafter\gdef\csname v:Mac:omp:e2e:4096:1\endcsname{139.0}
\expandafter\gdef\csname v:Mac:omp:e2e:4096:10\endcsname{24.92}
\expandafter\gdef\csname v:Mac:omp:e2e:4096:2\endcsname{71.53}
\expandafter\gdef\csname v:Mac:omp:e2e:4096:4\endcsname{36.00}
\expandafter\gdef\csname v:Mac:omp:e2e:4096:6\endcsname{34.80}
\expandafter\gdef\csname v:Mac:omp:e2e:4096:8\endcsname{26.72}
\expandafter\gdef\csname v:Mac:omp:e2e:64:1\endcsname{139.0}
\expandafter\gdef\csname v:Mac:omp:e2e:64:10\endcsname{47.70}
\expandafter\gdef\csname v:Mac:omp:e2e:64:2\endcsname{74.94}
\expandafter\gdef\csname v:Mac:omp:e2e:64:4\endcsname{42.01}
\expandafter\gdef\csname v:Mac:omp:e2e:64:6\endcsname{51.33}
\expandafter\gdef\csname v:Mac:omp:e2e:64:8\endcsname{45.22}
\expandafter\gdef\csname v:Mac:omp:e2e:8:1\endcsname{139.8}
\expandafter\gdef\csname v:Mac:omp:e2e:8:10\endcsname{155.5}
\expandafter\gdef\csname v:Mac:omp:e2e:8:2\endcsname{101.4}
\expandafter\gdef\csname v:Mac:omp:e2e:8:4\endcsname{80.51}
\expandafter\gdef\csname v:Mac:omp:e2e:8:6\endcsname{118.9}
\expandafter\gdef\csname v:Mac:omp:e2e:8:8\endcsname{134.6}
\expandafter\gdef\csname v:Mac:omp:eff:1024:1\endcsname{1.00}
\expandafter\gdef\csname v:Mac:omp:eff:1024:10\endcsname{0.54}
\expandafter\gdef\csname v:Mac:omp:eff:1024:2\endcsname{0.98}
\expandafter\gdef\csname v:Mac:omp:eff:1024:4\endcsname{0.96}
\expandafter\gdef\csname v:Mac:omp:eff:1024:6\endcsname{0.65}
\expandafter\gdef\csname v:Mac:omp:eff:1024:8\endcsname{0.63}
\expandafter\gdef\csname v:Mac:omp:eff:128:1\endcsname{1.00}
\expandafter\gdef\csname v:Mac:omp:eff:128:10\endcsname{0.38}
\expandafter\gdef\csname v:Mac:omp:eff:128:2\endcsname{0.95}
\expandafter\gdef\csname v:Mac:omp:eff:128:4\endcsname{0.89}
\expandafter\gdef\csname v:Mac:omp:eff:128:6\endcsname{0.53}
\expandafter\gdef\csname v:Mac:omp:eff:128:8\endcsname{0.48}
\expandafter\gdef\csname v:Mac:omp:eff:1:1\endcsname{1.04}
\expandafter\gdef\csname v:Mac:omp:eff:1:10\endcsname{0.01}
\expandafter\gdef\csname v:Mac:omp:eff:1:2\endcsname{0.27}
\expandafter\gdef\csname v:Mac:omp:eff:1:4\endcsname{0.10}
\expandafter\gdef\csname v:Mac:omp:eff:1:6\endcsname{0.04}
\expandafter\gdef\csname v:Mac:omp:eff:1:8\endcsname{0.02}
\expandafter\gdef\csname v:Mac:omp:eff:256:1\endcsname{1.00}
\expandafter\gdef\csname v:Mac:omp:eff:256:10\endcsname{0.44}
\expandafter\gdef\csname v:Mac:omp:eff:256:2\endcsname{0.97}
\expandafter\gdef\csname v:Mac:omp:eff:256:4\endcsname{0.93}
\expandafter\gdef\csname v:Mac:omp:eff:256:6\endcsname{0.61}
\expandafter\gdef\csname v:Mac:omp:eff:256:8\endcsname{0.54}
\expandafter\gdef\csname v:Mac:omp:eff:32:1\endcsname{1.00}
\expandafter\gdef\csname v:Mac:omp:eff:32:10\endcsname{0.22}
\expandafter\gdef\csname v:Mac:omp:eff:32:2\endcsname{0.88}
\expandafter\gdef\csname v:Mac:omp:eff:32:4\endcsname{0.74}
\expandafter\gdef\csname v:Mac:omp:eff:32:6\endcsname{0.36}
\expandafter\gdef\csname v:Mac:omp:eff:32:8\endcsname{0.30}
\expandafter\gdef\csname v:Mac:omp:eff:4096:1\endcsname{1.00}
\expandafter\gdef\csname v:Mac:omp:eff:4096:10\endcsname{0.56}
\expandafter\gdef\csname v:Mac:omp:eff:4096:2\endcsname{0.97}
\expandafter\gdef\csname v:Mac:omp:eff:4096:4\endcsname{0.96}
\expandafter\gdef\csname v:Mac:omp:eff:4096:6\endcsname{0.66}
\expandafter\gdef\csname v:Mac:omp:eff:4096:8\endcsname{0.65}
\expandafter\gdef\csname v:Mac:omp:eff:64:1\endcsname{1.00}
\expandafter\gdef\csname v:Mac:omp:eff:64:10\endcsname{0.29}
\expandafter\gdef\csname v:Mac:omp:eff:64:2\endcsname{0.93}
\expandafter\gdef\csname v:Mac:omp:eff:64:4\endcsname{0.83}
\expandafter\gdef\csname v:Mac:omp:eff:64:6\endcsname{0.45}
\expandafter\gdef\csname v:Mac:omp:eff:64:8\endcsname{0.38}
\expandafter\gdef\csname v:Mac:omp:eff:8:1\endcsname{0.99}
\expandafter\gdef\csname v:Mac:omp:eff:8:10\endcsname{0.09}
\expandafter\gdef\csname v:Mac:omp:eff:8:2\endcsname{0.68}
\expandafter\gdef\csname v:Mac:omp:eff:8:4\endcsname{0.43}
\expandafter\gdef\csname v:Mac:omp:eff:8:6\endcsname{0.19}
\expandafter\gdef\csname v:Mac:omp:eff:8:8\endcsname{0.13}
\expandafter\gdef\csname v:Mac:omp:speedup:1024:1\endcsname{1.00}
\expandafter\gdef\csname v:Mac:omp:speedup:1024:10\endcsname{5.40}
\expandafter\gdef\csname v:Mac:omp:speedup:1024:2\endcsname{1.95}
\expandafter\gdef\csname v:Mac:omp:speedup:1024:4\endcsname{3.84}
\expandafter\gdef\csname v:Mac:omp:speedup:1024:6\endcsname{3.91}
\expandafter\gdef\csname v:Mac:omp:speedup:1024:8\endcsname{5.04}
\expandafter\gdef\csname v:Mac:omp:speedup:128:1\endcsname{1.00}
\expandafter\gdef\csname v:Mac:omp:speedup:128:10\endcsname{3.83}
\expandafter\gdef\csname v:Mac:omp:speedup:128:2\endcsname{1.91}
\expandafter\gdef\csname v:Mac:omp:speedup:128:4\endcsname{3.55}
\expandafter\gdef\csname v:Mac:omp:speedup:128:6\endcsname{3.18}
\expandafter\gdef\csname v:Mac:omp:speedup:128:8\endcsname{3.82}
\expandafter\gdef\csname v:Mac:omp:speedup:1:1\endcsname{1.04}
\expandafter\gdef\csname v:Mac:omp:speedup:1:10\endcsname{0.13}
\expandafter\gdef\csname v:Mac:omp:speedup:1:2\endcsname{0.53}
\expandafter\gdef\csname v:Mac:omp:speedup:1:4\endcsname{0.40}
\expandafter\gdef\csname v:Mac:omp:speedup:1:6\endcsname{0.23}
\expandafter\gdef\csname v:Mac:omp:speedup:1:8\endcsname{0.16}
\expandafter\gdef\csname v:Mac:omp:speedup:256:1\endcsname{1.00}
\expandafter\gdef\csname v:Mac:omp:speedup:256:10\endcsname{4.37}
\expandafter\gdef\csname v:Mac:omp:speedup:256:2\endcsname{1.94}
\expandafter\gdef\csname v:Mac:omp:speedup:256:4\endcsname{3.70}
\expandafter\gdef\csname v:Mac:omp:speedup:256:6\endcsname{3.65}
\expandafter\gdef\csname v:Mac:omp:speedup:256:8\endcsname{4.32}
\expandafter\gdef\csname v:Mac:omp:speedup:32:1\endcsname{1.00}
\expandafter\gdef\csname v:Mac:omp:speedup:32:10\endcsname{2.18}
\expandafter\gdef\csname v:Mac:omp:speedup:32:2\endcsname{1.75}
\expandafter\gdef\csname v:Mac:omp:speedup:32:4\endcsname{2.95}
\expandafter\gdef\csname v:Mac:omp:speedup:32:6\endcsname{2.19}
\expandafter\gdef\csname v:Mac:omp:speedup:32:8\endcsname{2.44}
\expandafter\gdef\csname v:Mac:omp:speedup:4096:1\endcsname{1.00}
\expandafter\gdef\csname v:Mac:omp:speedup:4096:10\endcsname{5.56}
\expandafter\gdef\csname v:Mac:omp:speedup:4096:2\endcsname{1.94}
\expandafter\gdef\csname v:Mac:omp:speedup:4096:4\endcsname{3.85}
\expandafter\gdef\csname v:Mac:omp:speedup:4096:6\endcsname{3.98}
\expandafter\gdef\csname v:Mac:omp:speedup:4096:8\endcsname{5.19}
\expandafter\gdef\csname v:Mac:omp:speedup:64:1\endcsname{1.00}
\expandafter\gdef\csname v:Mac:omp:speedup:64:10\endcsname{2.91}
\expandafter\gdef\csname v:Mac:omp:speedup:64:2\endcsname{1.85}
\expandafter\gdef\csname v:Mac:omp:speedup:64:4\endcsname{3.30}
\expandafter\gdef\csname v:Mac:omp:speedup:64:6\endcsname{2.70}
\expandafter\gdef\csname v:Mac:omp:speedup:64:8\endcsname{3.07}
\expandafter\gdef\csname v:Mac:omp:speedup:8:1\endcsname{0.99}
\expandafter\gdef\csname v:Mac:omp:speedup:8:10\endcsname{0.89}
\expandafter\gdef\csname v:Mac:omp:speedup:8:2\endcsname{1.37}
\expandafter\gdef\csname v:Mac:omp:speedup:8:4\endcsname{1.72}
\expandafter\gdef\csname v:Mac:omp:speedup:8:6\endcsname{1.17}
\expandafter\gdef\csname v:Mac:omp:speedup:8:8\endcsname{1.03}
\expandafter\gdef\csname v:Mac:runs\endcsname{20}
\expandafter\gdef\csname v:Mac:serial:e2e:1\endcsname{140.7}
\expandafter\gdef\csname v:Mac:serial:e2e:1024\endcsname{139.1}
\expandafter\gdef\csname v:Mac:serial:e2e:128\endcsname{138.6}
\expandafter\gdef\csname v:Mac:serial:e2e:256\endcsname{138.7}
\expandafter\gdef\csname v:Mac:serial:e2e:32\endcsname{139.5}
\expandafter\gdef\csname v:Mac:serial:e2e:4096\endcsname{138.7}
\expandafter\gdef\csname v:Mac:serial:e2e:64\endcsname{138.8}
\expandafter\gdef\csname v:Mac:serial:e2e:8\endcsname{138.8}
\expandafter\gdef\csname v:Mac:serial:l1:1\endcsname{122.5}
\expandafter\gdef\csname v:Mac:serial:l1:1024\endcsname{120.6}
\expandafter\gdef\csname v:Mac:serial:l1:128\endcsname{120.2}
\expandafter\gdef\csname v:Mac:serial:l1:256\endcsname{120.3}
\expandafter\gdef\csname v:Mac:serial:l1:32\endcsname{121.0}
\expandafter\gdef\csname v:Mac:serial:l1:4096\endcsname{120.3}
\expandafter\gdef\csname v:Mac:serial:l1:64\endcsname{120.4}
\expandafter\gdef\csname v:Mac:serial:l1:8\endcsname{120.4}
\expandafter\gdef\csname v:Mac:serial:l1share:1\endcsname{87.1}
\expandafter\gdef\csname v:Mac:serial:l1share:1024\endcsname{86.7}
\expandafter\gdef\csname v:Mac:serial:l1share:128\endcsname{86.7}
\expandafter\gdef\csname v:Mac:serial:l1share:256\endcsname{86.7}
\expandafter\gdef\csname v:Mac:serial:l1share:32\endcsname{86.7}
\expandafter\gdef\csname v:Mac:serial:l1share:4096\endcsname{86.7}
\expandafter\gdef\csname v:Mac:serial:l1share:64\endcsname{86.7}
\expandafter\gdef\csname v:Mac:serial:l1share:8\endcsname{86.7}
\expandafter\gdef\csname v:Mac:serial:l2:1\endcsname{9.22}
\expandafter\gdef\csname v:Mac:serial:l2:1024\endcsname{8.86}
\expandafter\gdef\csname v:Mac:serial:l2:128\endcsname{8.82}
\expandafter\gdef\csname v:Mac:serial:l2:256\endcsname{8.84}
\expandafter\gdef\csname v:Mac:serial:l2:32\endcsname{8.88}
\expandafter\gdef\csname v:Mac:serial:l2:4096\endcsname{8.81}
\expandafter\gdef\csname v:Mac:serial:l2:64\endcsname{8.84}
\expandafter\gdef\csname v:Mac:serial:l2:8\endcsname{8.86}
\expandafter\gdef\csname v:Mac:serial:pre:1\endcsname{8.85}
\expandafter\gdef\csname v:Mac:serial:pre:1024\endcsname{9.62}
\expandafter\gdef\csname v:Mac:serial:pre:128\endcsname{9.58}
\expandafter\gdef\csname v:Mac:serial:pre:256\endcsname{9.60}
\expandafter\gdef\csname v:Mac:serial:pre:32\endcsname{9.62}
\expandafter\gdef\csname v:Mac:serial:pre:4096\endcsname{9.58}
\expandafter\gdef\csname v:Mac:serial:pre:64\endcsname{9.60}
\expandafter\gdef\csname v:Mac:serial:pre:8\endcsname{9.54}
\expandafter\gdef\csname v:Mac:serial:preshare:1\endcsname{6.3}
\expandafter\gdef\csname v:Mac:serial:preshare:1024\endcsname{6.9}
\expandafter\gdef\csname v:Mac:serial:preshare:128\endcsname{6.9}
\expandafter\gdef\csname v:Mac:serial:preshare:256\endcsname{6.9}
\expandafter\gdef\csname v:Mac:serial:preshare:32\endcsname{6.9}
\expandafter\gdef\csname v:Mac:serial:preshare:4096\endcsname{6.9}
\expandafter\gdef\csname v:Mac:serial:preshare:64\endcsname{6.9}
\expandafter\gdef\csname v:Mac:serial:preshare:8\endcsname{6.9}
\expandafter\gdef\csname v:Mac:serial:tput:1\endcsname{71{,}054}
\expandafter\gdef\csname v:Mac:serial:tput:1024\endcsname{71{,}893}
\expandafter\gdef\csname v:Mac:serial:tput:128\endcsname{72{,}158}
\expandafter\gdef\csname v:Mac:serial:tput:256\endcsname{72{,}092}
\expandafter\gdef\csname v:Mac:serial:tput:32\endcsname{71{,}683}
\expandafter\gdef\csname v:Mac:serial:tput:4096\endcsname{72{,}107}
\expandafter\gdef\csname v:Mac:serial:tput:64\endcsname{72{,}043}
\expandafter\gdef\csname v:Mac:serial:tput:8\endcsname{72{,}049}
\expandafter\gdef\csname v:Mac:val:batch:agree\endcsname{100.00}
\expandafter\gdef\csname v:Mac:val:batch:maxerr\endcsname{0}
\expandafter\gdef\csname v:Mac:val:batch:meanerr\endcsname{0}
\expandafter\gdef\csname v:Mac:val:batch:n\endcsname{3}
\expandafter\gdef\csname v:Mac:val:numpy:agree\endcsname{100.00}
\expandafter\gdef\csname v:Mac:val:numpy:maxerr\endcsname{\ensuremath{9.2\times10^{-5}}}
\expandafter\gdef\csname v:Mac:val:numpy:meanerr\endcsname{\ensuremath{1.7\times10^{-5}}}
\expandafter\gdef\csname v:Mac:val:numpy:n\endcsname{1}
\expandafter\gdef\csname v:Mac:val:omp:agree\endcsname{100.00}
\expandafter\gdef\csname v:Mac:val:omp:maxerr\endcsname{0}
\expandafter\gdef\csname v:Mac:val:omp:meanerr\endcsname{0}
\expandafter\gdef\csname v:Mac:val:omp:n\endcsname{16}
\expandafter\gdef\csname v:Mac:val:passed\endcsname{20}
\expandafter\gdef\csname v:Mac:val:total\endcsname{20}
\expandafter\gdef\csname v:Sp:val:batch:agree\endcsname{100.00}
\expandafter\gdef\csname v:Sp:val:batch:maxerr\endcsname{0}
\expandafter\gdef\csname v:Sp:val:batch:meanerr\endcsname{0}
\expandafter\gdef\csname v:Sp:val:batch:n\endcsname{3}
\expandafter\gdef\csname v:Sp:val:cuda:agree\endcsname{100.00}
\expandafter\gdef\csname v:Sp:val:cuda:maxerr\endcsname{\ensuremath{7.6\times10^{-6}}}
\expandafter\gdef\csname v:Sp:val:cuda:meanerr\endcsname{\ensuremath{5.0\times10^{-7}}}
\expandafter\gdef\csname v:Sp:val:cuda:n\endcsname{12}
\expandafter\gdef\csname v:Sp:val:hseq:agree\endcsname{100.00}
\expandafter\gdef\csname v:Sp:val:hseq:maxerr\endcsname{\ensuremath{7.6\times10^{-6}}}
\expandafter\gdef\csname v:Sp:val:hseq:meanerr\endcsname{\ensuremath{5.0\times10^{-7}}}
\expandafter\gdef\csname v:Sp:val:hseq:n\endcsname{16}
\expandafter\gdef\csname v:Sp:val:hyb:agree\endcsname{100.00}
\expandafter\gdef\csname v:Sp:val:hyb:maxerr\endcsname{\ensuremath{7.6\times10^{-6}}}
\expandafter\gdef\csname v:Sp:val:hyb:meanerr\endcsname{\ensuremath{5.0\times10^{-7}}}
\expandafter\gdef\csname v:Sp:val:hyb:n\endcsname{32}
\expandafter\gdef\csname v:Sp:val:numpy:agree\endcsname{100.00}
\expandafter\gdef\csname v:Sp:val:numpy:maxerr\endcsname{\ensuremath{9.3\times10^{-5}}}
\expandafter\gdef\csname v:Sp:val:numpy:meanerr\endcsname{\ensuremath{1.7\times10^{-5}}}
\expandafter\gdef\csname v:Sp:val:numpy:n\endcsname{1}
\expandafter\gdef\csname v:Sp:val:omp:agree\endcsname{100.00}
\expandafter\gdef\csname v:Sp:val:omp:maxerr\endcsname{0}
\expandafter\gdef\csname v:Sp:val:omp:meanerr\endcsname{0}
\expandafter\gdef\csname v:Sp:val:omp:n\endcsname{16}
\expandafter\gdef\csname v:Sp:val:passed\endcsname{80}
\expandafter\gdef\csname v:Sp:val:total\endcsname{80}
\expandafter\gdef\csname v:SpO3:val:batch:agree\endcsname{100.00}
\expandafter\gdef\csname v:SpO3:val:batch:maxerr\endcsname{0}
\expandafter\gdef\csname v:SpO3:val:batch:meanerr\endcsname{0}
\expandafter\gdef\csname v:SpO3:val:batch:n\endcsname{3}
\expandafter\gdef\csname v:SpO3:val:cuda:agree\endcsname{100.00}
\expandafter\gdef\csname v:SpO3:val:cuda:maxerr\endcsname{\ensuremath{1.1\times10^{-5}}}
\expandafter\gdef\csname v:SpO3:val:cuda:meanerr\endcsname{\ensuremath{9.9\times10^{-7}}}
\expandafter\gdef\csname v:SpO3:val:cuda:n\endcsname{12}
\expandafter\gdef\csname v:SpO3:val:hseq:agree\endcsname{100.00}
\expandafter\gdef\csname v:SpO3:val:hseq:maxerr\endcsname{\ensuremath{1.1\times10^{-5}}}
\expandafter\gdef\csname v:SpO3:val:hseq:meanerr\endcsname{\ensuremath{9.9\times10^{-7}}}
\expandafter\gdef\csname v:SpO3:val:hseq:n\endcsname{16}
\expandafter\gdef\csname v:SpO3:val:hyb:agree\endcsname{100.00}
\expandafter\gdef\csname v:SpO3:val:hyb:maxerr\endcsname{\ensuremath{1.1\times10^{-5}}}
\expandafter\gdef\csname v:SpO3:val:hyb:meanerr\endcsname{\ensuremath{9.9\times10^{-7}}}
\expandafter\gdef\csname v:SpO3:val:hyb:n\endcsname{32}
\expandafter\gdef\csname v:SpO3:val:numpy:agree\endcsname{100.00}
\expandafter\gdef\csname v:SpO3:val:numpy:maxerr\endcsname{\ensuremath{9.3\times10^{-5}}}
\expandafter\gdef\csname v:SpO3:val:numpy:meanerr\endcsname{\ensuremath{1.7\times10^{-5}}}
\expandafter\gdef\csname v:SpO3:val:numpy:n\endcsname{1}
\expandafter\gdef\csname v:SpO3:val:omp:agree\endcsname{100.00}
\expandafter\gdef\csname v:SpO3:val:omp:maxerr\endcsname{0}
\expandafter\gdef\csname v:SpO3:val:omp:meanerr\endcsname{0}
\expandafter\gdef\csname v:SpO3:val:omp:n\endcsname{16}
\expandafter\gdef\csname v:SpO3:val:passed\endcsname{80}
\expandafter\gdef\csname v:SpO3:val:total\endcsname{80}

\maketitle

\begin{abstract}
Deep-learning inference can run on multicore CPUs, on GPUs, or on both, but GPU speedups are often quoted for the kernels alone. We compare serial C++, OpenMP, CUDA and a hybrid OpenMP--CUDA software pipeline on the same MNIST network (784$\to$256$\to$10, 10{,}000 images, batches of 1 to 4096). Each configuration is measured \val{Sp:runs} times on a node with a Tesla K40m GPU (CUDA 11.7) and 24 CPU threads, and the CPU codes are also run on an Apple M4. All implementations agree with the serial reference (largest logit difference \val{Sp:val:cuda:maxerr}, identical predictions). With 24 threads OpenMP is \val{Sp:omp:speedup:4096:24}$\times$ faster than serial at batch 4096 (efficiency \val{Sp:omp:eff:4096:24}). CUDA is \val{Sp:cuda:speedup:4096}$\times$ faster end-to-end, although its kernels alone are \val{Sp:cuda:kspeedup:4096}$\times$ faster, and for single images it is slower than serial. Copies take at most \val{Sp:cuda:tshare:1}\,\% of the CUDA time, whereas single-threaded preprocessing grows to \val{Sp:cuda:preshare:4096}\,\%. The pipelined hybrid needs \val{Sp:hyb:best:e2e:4096}\,ms per pass at batch 4096, \val{Sp:hyb:vscuda:4096}$\times$ faster than CUDA-only and \val{Sp:hyb:vscpu:4096}$\times$ faster than the best OpenMP configuration, but OpenMP is faster for batches of 1 and 8.

\end{abstract}

\section{Introduction}
Inference with deep neural networks is dominated by dense linear algebra, which maps well onto parallel hardware. A modern machine offers two kinds of parallelism: a multicore CPU, programmed here with the OpenMP shared-memory directives~\cite{dagum1998openmp}, and a GPU, programmed with CUDA~\cite{nickolls2008scalable}. A GPU runs thousands of threads, but its data must first cross the PCIe bus, and each kernel launch has a fixed cost. GPU speedups are often quoted for the kernel alone or against a weakly optimized CPU baseline, which can overstate the benefit~\cite{lee2010debunking,gregg2011data}. Heterogeneous execution tries to keep both processors busy~\cite{mittal2015survey}, but whether it pays off depends on how the work is split and overlapped.

This paper measures these effects on a small but representative workload: inference of a fully connected network (784$\to$256$\to$10) on the 10{,}000 MNIST test images, processed in batches of 1 to 4096. We implement the same computation four ways---serial C++, OpenMP, CUDA, and a hybrid OpenMP--CUDA software pipeline---and compare them on identical data and weights. We ask six research questions:

\begin{itemize}
\item \textbf{RQ1} How much speedup does OpenMP give over serial execution?
\item \textbf{RQ2} How does CUDA compare with serial and OpenMP execution, for kernel-only and end-to-end time?
\item \textbf{RQ3} How does the batch size affect CPU and GPU performance?
\item \textbf{RQ4} How much of the GPU execution time is data transfer rather than computation?
\item \textbf{RQ5} Does a hybrid OpenMP--CUDA implementation improve end-to-end performance?
\item \textbf{RQ6} How do the results relate to Amdahl's law, parallel overhead, and transfer costs?
\end{itemize}

\noindent Our contributions are: (i)~four implementations of one workload whose outputs are verified against each other and against an independent float64 computation; (ii)~a measurement protocol with warm-up, 20 repetitions and confidence intervals, following the recommendations of Hoefler and Belli~\cite{hoefler2015scientific}; (iii)~measured answers to RQ1--RQ6 on two machines; and (iv)~a reproducible code and analysis package. Every number in this paper is generated from the raw measurements by script.


\section{Background and Related Work}
\subsection{OpenMP and CUDA}
OpenMP is a directive-based API for shared-memory parallelism~\cite{dagum1998openmp,openmparb2021specification}. A \texttt{parallel} region forks a team of threads; a work-shared loop (\texttt{omp for}) divides its iterations among them and ends with an implicit barrier. Forking and barriers cost time, so the work per region must be large enough to amortize them.

CUDA organizes GPU work as a grid of thread blocks, each block a group of threads that run the same kernel on different data~\cite{nickolls2008scalable,nvidia2026cudaguide}. Host and device memories are separate, so inputs must be copied to the device and results copied back. Kernel launches and copies on non-default streams are asynchronous, but asynchronous copies need page-locked (pinned) host memory~\cite{nvidia2026cudaguide}; streams allow the copy of one batch to overlap with the kernels of another~\cite{harris2012overlap}. Pinned memory also speeds up transfers~\cite{harris2012transfers}. We time GPU stages with CUDA events~\cite{harris2012metrics}.

\subsection{Scaling Laws}
Amdahl's law~\cite{amdahl1967validity} bounds the speedup on $p$ processors by $S(p)=1/(f+(1-f)/p)$, where $f$ is the fraction of the work that stays serial; for a fixed problem size $S(p)<1/f$. Gustafson~\cite{gustafson1988reevaluating} observed that enlarging the problem reduces the effective serial fraction, which is what a larger batch does here. Karp and Flatt~\cite{karp1990measuring} proposed the experimentally determined serial fraction $e(p)=(1/S-1/p)/(1-1/p)$; if $e$ grows with $p$, overheads other than a fixed serial part are present. Hill and Marty~\cite{hill2008amdahl} extend Amdahl's law to multicore chips and conclude that optimal performance needs both more parallelism and faster sequential execution. The roofline model~\cite{williams2009roofline} relates performance to memory bandwidth and arithmetic intensity.

\subsection{Related Work}
Lee et al.~\cite{lee2010debunking} found that after optimizing both platforms the GPU advantage over a multicore CPU on throughput kernels was about 2.5$\times$ on average, far below the large speedups that had been claimed for GPUs; this motivates our use of an OpenMP baseline besides the serial one. Gregg and Hazelwood~\cite{gregg2011data} show that transfer time can erase kernel speedups, so we report kernel-only and end-to-end results side by side. Mittal and Vetter~\cite{mittal2015survey} survey CPU--GPU collaboration techniques such as workload partitioning and pipelining. Sze et al.~\cite{sze2017efficient} point out that fully connected layers have no weight reuse at batch size one, and Jouppi et al.~\cite{jouppi2017tpu} report that MLP inference is limited by memory bandwidth and by latency targets. Our study complements these works with a controlled comparison of four implementations on identical inputs, including the effect of overlapping preprocessing with GPU work.


\section{Workload and Dataset}
\textbf{Dataset.} MNIST~\cite{lecun1998gradient} contains 28$\times$28 grayscale images of handwritten digits (60{,}000 for training, 10{,}000 for testing). We run inference on the 10{,}000 test images. The files are downloaded with a script and verified by MD5 and SHA-256 checksums.

\textbf{Network.} A fully connected network 784$\to$256$\to$10 with a ReLU after the hidden layer and float32 arithmetic (Figure~\ref{fig:network}). It has 203{,}530 parameters (814{,}120 bytes). The weights were trained once, outside the timed code, with a NumPy script (no deep-learning framework): He initialisation~\cite{he2015delving}, Adam~\cite{kingma2015adam} with learning rate $10^{-3}$, batch size 128, 10 epochs and seed 0. These hyper-parameters were fixed in advance and not tuned on the test set. The trained network classifies \val{Mac:accuracy}\,\% of the test images correctly. All implementations load the same weight file.

\textbf{Pipeline.} A \emph{pass} processes all 10{,}000 images in consecutive batches of size $B$ and performs four stages: (0)~per-image standardisation of the raw 8-bit pixels, $x_i\leftarrow(p_i-\mu)/\sqrt{\sigma^2+\epsilon}$ with $\mu$ and $\sigma^2$ computed per image; (1)~the hidden layer $h=\mathrm{ReLU}(W_1^{\top}x+b_1)$; (2)~the output layer $z=W_2^{\top}h+b_2$; (3)~the prediction $\arg\max z$. We use batch sizes 1, 8, 32, 64, 128, 256, 1024 and 4096; when $B$ does not divide 10{,}000 the last batch is smaller.

\begin{coltable}{Work per image, from the layer dimensions.}{tab:work}
\begin{tabular}{@{}lrr@{}}
\toprule
Stage & Multiply-adds & Parameter bytes \\
\midrule
Hidden layer (784$\to$256) & 200{,}704 & 803{,}840 \\
Output layer (256$\to$10) & 2{,}560 & 10{,}280 \\
\midrule
Total & 203{,}264 & 814{,}120 \\
\bottomrule
\end{tabular}

\smallskip
\footnotesize Input: 784 B (uint8) before and 3{,}136 B (float32) after stage~0; output: 44 B (10 logits and the class). One pass: 4.07 GFLOP in stages 1--2, 31.4 MB copied to the GPU, 0.44 MB copied back.
\end{coltable}

\widefigure{network_architecture.pdf}{The workload: preprocessing and a 784$\to$256$\to$10 network with ReLU.}{fig:network}


\section{Implementations}
All four implementations share the arithmetic of stages 0--3 (header \texttt{model.hpp}): the hidden layer accumulates over $k=0,\dots,783$ and the output layer over $j=0,\dots,255$ in the same order everywhere. They therefore perform the same mathematical work, and any difference in results comes only from fused multiply-add contraction. Figure~\ref{fig:arch} shows which processor executes which stage.

\widefigure{architecture.pdf}{Execution architecture of the four implementations. Colour shows the resource: one CPU thread, OpenMP threads, GPU threads, or a PCIe transfer.}{fig:arch}

\subsection{Serial CPU}
Plain C++ loops over the batch, one thread, compiled with the same optimisation flags as the OpenMP code. Stage~1 is organised in blocks of 16 consecutive hidden units (one 64-byte cache line of a row of $W_1$) so that the inner loop is contiguous and can be vectorised by the compiler. No library routines are used.

\subsection{OpenMP}
One parallel region covers the whole pass; inside it the batch loop is executed by every thread and each stage is a work-shared loop (Listing~\ref{lst:omp}). The implicit barrier after each loop is required because the next stage reads data written by other threads. Every iteration writes a distinct output, so there are no data races: stage~0 is owned per image, stage~1 per (image, block of 16 hidden units), stage~2 per (image, class), and stage~3 per image. Splitting stage~1 over (image, block) offers 16 independent work items even for a batch of one image. The batch loop itself is not parallelised because consecutive batches reuse the same buffers.

\begin{lstlisting}[language=C++,caption={OpenMP region (simplified).},label={lst:omp}]
#pragma omp parallel num_threads(T)
for (int s = 0; s < n; s += batch) {
  #pragma omp for schedule(static)
  for (b = 0; b < B; ++b) preprocess(b);
  #pragma omp for collapse(2)
  for (b = 0; b < B; ++b)
    for (jb = 0; jb < 16; ++jb) hidden(b, jb);
  #pragma omp for collapse(2)
  for (b = 0; b < B; ++b)
    for (c = 0; c < 10; ++c) output(b, c);
  #pragma omp for
  for (b = 0; b < B; ++b) argmax(b);
}
\end{lstlisting}

\subsection{CUDA}
Preprocessing runs on one CPU thread, the batch is copied to the GPU with a synchronous copy from pageable memory, three kernels compute the result, and the logits and classes are copied back (Figure~\ref{fig:arch}). In the hidden-layer kernel a 2-D grid of $d\times d$ thread blocks covers the output matrix: thread $(t_x,t_y)$ of block $(b_x,b_y)$ computes one element, hidden unit $j=b_xd+t_x$ of sample $b=b_yd+t_y$ (Listing~\ref{lst:cuda}). Because $j$ varies fastest, consecutive threads read consecutive addresses of $W_1$, so the loads are coalesced. The output-layer kernel uses one thread per (sample, class) and the argmax kernel one thread per sample, both with 128 threads per block. Every thread writes a distinct element, so no atomics are needed. Every CUDA call is checked for errors. The block size $d$ is 8, 16 or 32; 16 is the default.

\begin{lstlisting}[language=CUDA,caption={Hidden-layer kernel.},label={lst:cuda}]
__global__ void layer1_kernel(const float* x,
    const float* w1, const float* b1, float* h, int B) {
  int j = blockIdx.x*blockDim.x + threadIdx.x; // unit
  int b = blockIdx.y*blockDim.y + threadIdx.y; // sample
  if (b >= B || j >= 256) return;
  float acc = b1[j];
  for (int k = 0; k < 784; ++k)
    acc += x[b*784 + k] * w1[k*256 + j];
  h[b*256 + j] = acc > 0.f ? acc : 0.f;  // ReLU
}
\end{lstlisting}

\subsection{Hybrid OpenMP--CUDA}
The hybrid assigns each processor the stage it is best placed to do. OpenMP threads preprocess the batch into a pinned host buffer (stage~0, the only CPU stage); the DMA engines copy it to the GPU and copy the results back asynchronously; CUDA threads run the three kernels. Batches are assigned round-robin to $S$ slots (default $S=3$), each with its own CUDA stream, device buffers and pinned buffers. While the GPU works on earlier batches, the CPU already preprocesses the next one (a software pipeline~\cite{harris2012overlap}); a slot is reused only after its previous batch has finished, which is checked with a CUDA event. To separate the two ingredients we also measure a \emph{hybrid-sequential} variant: OpenMP preprocessing, but synchronous pageable copies and no overlap. Comparing CUDA-only, hybrid-sequential and hybrid-pipelined thus isolates the effect of parallel preprocessing from the effect of overlap and pinned memory.


\section{Experimental Method}\label{sec:method}
\textbf{Platforms.} Table~\ref{tab:platforms} lists the two machines. The \emph{Spartan GPU node} is the reference platform: it has both a multicore CPU and a GPU, so serial, OpenMP, CUDA and hybrid runs are compared on the same machine, one benchmark job per node under Slurm~\cite{yoo2003slurm}. The \emph{Apple M4} laptop has no NVIDIA GPU; it is used only for the CPU implementations and is never mixed with GPU-node numbers. The M4 has 4 performance and 6 efficiency cores, so its thread-scaling curve is not that of a homogeneous machine.

\begin{widetable}{The two machines. ``Physical cores / logical threads'' are the hardware counts; ``OpenMP threads used'' is what the benchmark could use. The Spartan job was allocated 24 of the node's 28 CPUs.}{tab:platforms}
\begin{tabular*}{\textwidth}{@{\extracolsep{\fill}}p{0.17\textwidth}p{0.36\textwidth}p{0.39\textwidth}@{}}
\toprule
 & \textbf{Spartan GPU node (reference)} & \textbf{Apple M4 laptop (CPU only)} \\
\midrule
CPU & 2$\times$ Intel Xeon E5-2680 v4, 2.40\,GHz & Apple M4 \\
Physical cores / logical threads & 28 / 28 (14 per socket, 1 thread per core) & 10 / 10 (4 performance + 6 efficiency) \\
OpenMP threads used & 1, 2, 4, 8, 12, 16, 24 & 1, 2, 4, 6, 8, 10 \\
Cache & 32\,KiB L1d and 256\,KiB L2 per core; 35\,MiB L3 per socket; 2 NUMA nodes & 16\,MiB L2 (performance cluster), 4\,MiB L2 (efficiency cluster) \\
Memory & 125\,GiB & 16\,GiB \\
GPU & Tesla K40m: compute capability 3.5, 15 SMs, 11{,}441\,MiB, 384-bit bus, memory clock 3{,}004\,MHz & none \\
Operating system & Linux 5.14 (Slurm, partition \texttt{gpus}) & macOS 26.6.2 \\
CUDA & \texttt{nvcc} 11.7.99, runtime 11.7, driver 470.256.02 (CUDA 11.4); \texttt{-arch=sm\_35} & not available \\
Host compiler & g++ 11.4.1, OpenMP 4.5; \texttt{-O3 -std=c++14}, plus \texttt{-march=native} for the reference runs & Apple clang 21.0.0 with libomp 23.1.2; \texttt{-O3 -std=c++14} \\
\bottomrule
\end{tabular*}
\end{widetable}


\textbf{Protocol.} For every configuration (implementation, batch size, thread count, block size, number of streams) we run 3 unrecorded warm-up passes and then \val{Sp:runs} recorded passes, and store every pass as one row of a CSV file together with the stage times, the accuracy of that pass, and the hardware and compiler description~\cite{hoefler2015scientific}. Raw files are never overwritten. A pass processes the 10{,}000 test images; the timers start after all host buffers have been allocated, and one-time set-up (device buffers and the weight upload) is excluded for every implementation. OpenMP threads are bound with \texttt{OMP\_PROC\_BIND=close} and \texttt{OMP\_PLACES=cores} on Linux; macOS ignores thread binding.

\textbf{Timing.} \emph{End-to-end} time is the wall clock of a whole pass. CPU stages are timed with the steady clock; the GPU stages (host-to-device copy, the three kernels, device-to-host copy) are timed with CUDA events~\cite{harris2012metrics}. Kernel-only time is always reported together with end-to-end time, because they answer different questions: the first measures how fast the GPU computes, the second what an application actually waits for~\cite{gregg2011data}. In the pipelined hybrid the GPU stages overlap each other and the CPU work, so their sum can exceed the end-to-end time.

\textbf{Metrics.} With $T$ the end-to-end time of a pass: speedup $S=T_{\mathrm{serial}}/T_{\mathrm{impl}}$ on the same machine and batch size; OpenMP efficiency $E=S/p$ for $p$ threads; throughput $=10{,}000/T$ images per second; transfer share $=(T_{\mathrm{H2D}}+T_{\mathrm{D2H}})/T$. Each mean is given with a 95\,\% Student-$t$ confidence interval. Hybrid variants are compared with Welch's $t$-test and a bootstrap interval of the ratio of means (5{,}000 resamples). The Amdahl serial fraction $f$ is fitted by least squares to the measured speedups of each batch size, and the Karp--Flatt metric is computed for every thread count. When we say ``best'' thread count we mean the tested count with the lowest mean time.

\textbf{Compiler flags.} All CPU code is compiled with \texttt{-O3}. On the Spartan node, whose default x86-64 target only uses 128-bit vector instructions, the reference runs use \texttt{-march=native} so that the compiler may use the vector instructions of the host CPU for these straightforward loops~\cite{lee2010debunking}; a second set of runs with plain \texttt{-O3} quantifies how much this matters. No implementation uses a tuned BLAS routine or a hand-tiled kernel.


\section{Correctness Validation}
Performance numbers are only meaningful if every implementation computes the same result, so the correctness gate runs before any timing. It compares every implementation with the serial reference on all 10{,}000 test images. The acceptance criteria were fixed before the first measurement: a maximum absolute logit difference of at most $10^{-4}$ and a prediction agreement of at least 99.9\,\%. We test batch sizes 1, 37 (which does not divide 10{,}000, so the last batch is partial), 256 and 4096; OpenMP with 1, 2, 4 and all threads; CUDA with block sizes 8, 16 and 32; and both hybrid variants with 1 and 3 streams (pipelined only) and four thread counts. An independent anchor compares the serial float32 logits of the first 100 images with a float64 NumPy computation of the same network.

On the Spartan node all \val{Sp:val:total} comparisons pass (Table~\ref{tab:correct}). OpenMP is bit-identical to serial (maximum difference \val{Sp:val:omp:maxerr}) because every output element is accumulated in the same source-level order and both are compiled with the same flags. CUDA and both hybrid variants differ from serial by at most \val{Sp:val:cuda:maxerr} (mean \val{Sp:val:cuda:meanerr}), which is consistent with compiler-specific floating-point optimisations such as fused multiply-add contraction; in the plain \texttt{-O3} build, whose host code targets baseline x86-64, the largest difference is \val{SpO3:val:cuda:maxerr}. We did not set the contraction options explicitly (Table~\ref{tab:compiler}). The prediction agreement is \val{Sp:val:cuda:agree}\,\% everywhere, so the GPU and hybrid versions predict exactly the same digit as the serial code for every image, and the accuracy is \val{Sp:accuracy}\,\% in all runs (Figure~\ref{fig:correct}). The float32 implementation differs from the float64 anchor by at most \val{Sp:val:numpy:maxerr}, which sets the scale of float32 rounding for this network and is well inside the tolerance. On the Mac, the serial and OpenMP implementations also pass all \val{Mac:val:total} comparisons.

\begin{widetable}{Correctness gate on the Spartan node (reference: serial). ``Configs'' counts the batch-size, thread, block and stream combinations tested; errors are in logit units.}{tab:correct}
\begin{tabular}{lrrrrrr}
\toprule
Comparison & Configs & Max abs err & Mean abs err & Agree \% & Acc \% & Pass \\
\midrule
Serial vs NumPy (float64) & 1 & 9.3e-5 & 1.7e-5 & 100 & -- & 1/1 \\
Serial batch-invariance & 3 & 0 & 0 & 100 & 97.97 & 3/3 \\
Serial vs OpenMP & 16 & 0 & 0 & 100 & 97.97 & 16/16 \\
Serial vs CUDA & 12 & 7.6e-6 & 5.0e-7 & 100 & 97.97 & 12/12 \\
Serial vs Hybrid-sequential & 16 & 7.6e-6 & 5.0e-7 & 100 & 97.97 & 16/16 \\
Serial vs Hybrid-pipelined & 32 & 7.6e-6 & 5.0e-7 & 100 & 97.97 & 32/32 \\
\bottomrule
\end{tabular}

\end{widetable}

\colfigure{correctness_all.pdf}{Worst-case maximum logit difference against the serial reference, over all tested configurations on the Mac and the Spartan node.}{fig:correct}


\section{Performance Results}
All times are the mean time of one pass over 10{,}000 images, in milliseconds (ms), from \val{Sp:runs} recorded passes per configuration on the Spartan node unless the Apple M4 is named. Batch sizes are in images per batch; ``$S$'' is the speedup over the serial code on the same machine and batch size (in multiples, $\times$), and ``$E$'' is the efficiency $S/p$. Compute-only comparisons (stages 1--3, no preprocessing and no copies) and end-to-end comparisons (stages 0--3 including copies) are reported separately. The M4 is used only for the OpenMP study in Section~\ref{sec:omp}; no CPU-versus-GPU comparison uses it.

\subsection{OpenMP Scalability (RQ1)}\label{sec:omp}

On the Xeon the speedup is close to linear for batches of 32 images or more (Figure~\ref{fig:omp}, Table~\ref{tab:omp}). With 24 threads the speedup is \val{Sp:omp:speedup:4096:24}$\times$ at batch 4096 (efficiency \val{Sp:omp:eff:4096:24}) and \val{Sp:omp:speedup:256:24}$\times$ at batch 256 (efficiency \val{Sp:omp:eff:256:24}). With one image per batch the best result is \val{Sp:omp:best:speedup:1}$\times$ with \val{Sp:omp:best:thr:1} threads, and the efficiency at 24 threads falls to \val{Sp:omp:eff:1:24}: a batch of one image offers only 16 work items per stage and each stage ends in a barrier. OpenMP with one thread needs \val{Sp:omp:e2e:256:1}\,ms against \val{Sp:serial:e2e:256}\,ms for the serial code (batch 256), so the OpenMP runtime adds no measurable cost.

The Apple M4 behaves differently. The speedup follows the ideal line up to its four performance cores (\val{Mac:omp:speedup:4096:4}$\times$ with 4 threads at batch 4096) and then flattens to \val{Mac:omp:speedup:4096:10}$\times$ with 10 threads (efficiency \val{Mac:omp:eff:4096:10}); the dotted line in Figure~\ref{fig:omp} marks the four performance cores. With one image per batch OpenMP is \emph{slower} than serial on the M4 for every thread count above one: \val{Mac:omp:speedup:1:2}$\times$ with 2 threads and \val{Mac:omp:speedup:1:10}$\times$ with 10 threads. The serial pass takes \val{Mac:serial:e2e:4096}\,ms on the M4 and \val{Sp:serial:e2e:4096}\,ms on one Xeon core (a CPU-to-CPU comparison; the build settings differ, Table~\ref{tab:compiler}).

\widefigure{omp_combined.pdf}{OpenMP speedup over serial (top) and parallel efficiency (bottom) versus thread count, for four batch sizes, on the Apple M4 (left) and the Xeon (right). The dashed lines are ideal scaling; the dotted line marks the four performance cores of the M4.}{fig:omp}

\begin{widetable}{OpenMP scaling on the Spartan node: time per pass (ms), speedup $S$ over serial ($\times$) and efficiency $E=S/p$ for three batch sizes.}{tab:omp}
\begin{tabular}{rrrrrrrrrr}
\toprule
Threads & \multicolumn{3}{c}{Batch 1} & \multicolumn{3}{c}{Batch 256} & \multicolumn{3}{c}{Batch 4096} \\
\cmidrule(lr){2-4}\cmidrule(lr){5-7}\cmidrule(lr){8-10}
 & Time (ms) & $S$ ($\times$) & $E$ & Time (ms) & $S$ ($\times$) & $E$ & Time (ms) & $S$ ($\times$) & $E$ \\
\midrule
1 & 1{,}221.8 & 1.01 & 1.01 & 1{,}205.5 & 1.02 & 1.02 & 1{,}205.5 & 1.02 & 1.02 \\
2 & 662.6 & 1.86 & 0.93 & 631.0 & 1.96 & 0.98 & 629.5 & 1.96 & 0.98 \\
4 & 370.8 & 3.33 & 0.83 & 333.8 & 3.70 & 0.93 & 334.0 & 3.70 & 0.92 \\
8 & 212.6 & 5.81 & 0.73 & 168.6 & 7.32 & 0.92 & 166.8 & 7.41 & 0.93 \\
12 & 216.1 & 5.72 & 0.48 & 114.7 & 10.77 & 0.90 & 112.6 & 10.98 & 0.91 \\
16 & 162.8 & 7.59 & 0.47 & 90.3 & 13.68 & 0.86 & 85.5 & 14.45 & 0.90 \\
24 & 167.8 & 7.36 & 0.31 & 58.9 & 20.99 & 0.87 & 57.3 & 21.55 & 0.90 \\
\bottomrule
\end{tabular}

\end{widetable}

\subsection{End-to-End Comparison}\label{sec:e2e}

This is the \emph{end-to-end} comparison of the four primary implementations (stages 0--3, copies included; the wall-clock time of a pass). Table~\ref{tab:cpugpu} and Figures~\ref{fig:exec}--\ref{fig:tput} show the results. The CUDA entries use a fixed 16$\times$16 block and one preprocessing thread, whereas OpenMP and the hybrid use their best thread count, so the comparison is ``CUDA with a fixed block against OpenMP with tuned thread count''. The serial and OpenMP times hardly depend on the batch size: the serial pass takes \val{Sp:serial:e2e:1}\,ms with batch 1 and \val{Sp:serial:e2e:4096}\,ms with batch 4096, because every image is processed independently by the same loops. The GPU depends strongly on the batch size: \cudabase{} falls from \val{Sp:cuda:e2e:1}\,ms (batch 1) to \val{Sp:cuda:e2e:4096}\,ms (batch 4096).

With one image per batch the GPU pipeline is slower than serial (\val{Sp:cuda:speedup:1}$\times$): each of the 10{,}000 batches needs three kernel launches, two copies and a synchronization, and the kernels alone take \val{Sp:cuda:kernel:1}\,ms per pass, \val{Sp:cuda:kus:1}\,$\mu$s per batch. From batch 8 on it is faster than serial, reaching \val{Sp:cuda:speedup:4096}$\times$ end-to-end at batch 4096.

Against the best OpenMP configuration, the CUDA pipeline is slower for batches up to 64, level at batch 128 (\val{Sp:cuda:e2e:128}\,ms against \val{Sp:omp:best:e2e:128}\,ms) and ahead from batch 256 (\val{Sp:cuda:e2e:256}\,ms against \val{Sp:omp:best:e2e:256}\,ms); at batch 4096 it is only \val{Sp:cuda:vscpu:4096}$\times$ faster than 24-thread OpenMP. Using the rate convention of Section~\ref{sec:method} (stages 1--2 only, divided by the end-to-end time unless marked kernel-only), a pass at batch 4096 reaches \val{Sp:serial:gflops:4096}\,GFLOP/s serial, \val{Sp:omp:best:gflops:4096}\,GFLOP/s with OpenMP and \val{Sp:cuda:gflops:4096}\,GFLOP/s for the CUDA pipeline end-to-end; the kernel-only rate of the CUDA kernels is \val{Sp:cuda:kgflops:4096}\,GFLOP/s.

Against the same baselines the pipelined hybrid is the fastest implementation from batch \val{Sp:hyb:crossover} on, \val{Sp:hyb:vscpu:4096}$\times$ faster than the best OpenMP configuration at batch 4096, whereas OpenMP is fastest for batches of 1 and 8 (Section~\ref{sec:hybrid}). Figure~\ref{fig:speedup} shows the speedup over serial of the four implementations and Figure~\ref{fig:tput} the throughput in images per second. At batch 4096 the serial code processes \val{Sp:serial:tput:4096}~images/s, OpenMP \val{Sp:omp:best:tput:4096}~images/s, the CUDA implementation with single-threaded preprocessing \val{Sp:cuda:tput:4096}~images/s and the pipelined hybrid \val{Sp:hyb:best:tput:4096}~images/s. Speedup and throughput are two views of the same time measurements: $S=T_{\mathrm{serial}}/T$ and throughput $=10{,}000/T$.

\begin{widetable}{End-to-end comparison on the Spartan node: mean time per pass (ms) and speedup over serial ($\times$). OpenMP and Hybrid show the best of the tested thread counts (in brackets). ``CUDA end-to-end'' preprocesses on one CPU thread; the kernel-only time of Table~\ref{tab:compute} excludes preprocessing and copies. $^{\ast}$The kernel-only speedup is a compute-only ratio (serial time of stages 1--3 divided by the kernel time); all other speedups are end-to-end ratios.}{tab:cpugpu}
\begin{tabular}{rrrrrrrrr}
\toprule
Batch & Serial & OpenMP (threads) & CUDA kernel & CUDA end-to-end & Hybrid (threads) & \multicolumn{3}{c}{Speedup over serial ($\times$)} \\
\cmidrule(lr){7-9}
 & (ms) & (ms) & only (ms) & (ms) & (ms) & kernel-only & CUDA e2e & Hybrid \\
\midrule
1 & 1{,}234.8 & 162.8 (16) & 1{,}091.0 & 1{,}507.0 & 650.9 (1) & 1.13 & 0.82 & 1.90 \\
8 & 1{,}235.7 & 75.5 (24) & 143.9 & 218.8 & 87.5 (8) & 8.59 & 5.65 & 14.13 \\
32 & 1{,}235.0 & 61.9 (24) & 47.3 & 85.5 & 33.3 (2) & 26.10 & 14.45 & 37.09 \\
64 & 1{,}234.8 & 59.2 (24) & 34.4 & 65.9 & 27.9 (24) & 35.87 & 18.73 & 44.20 \\
128 & 1{,}234.6 & 58.2 (24) & 29.8 & 58.6 & 26.5 (24) & 41.43 & 21.06 & 46.59 \\
256 & 1{,}235.2 & 58.9 (24) & 28.1 & 55.1 & 26.0 (24) & 44.01 & 22.40 & 47.54 \\
1024 & 1{,}235.0 & 57.7 (24) & 23.7 & 48.3 & 25.2 (24) & 52.02 & 25.56 & 48.98 \\
4096 & 1{,}235.4 & 57.3 (24) & 23.0 & 46.7 & 24.7 (24) & 53.65 & 26.44 & 50.02 \\
\bottomrule
\end{tabular}

\end{widetable}

\colfigure{exec_time_spartan_k40_native.pdf}{End-to-end time per pass (ms, log scale) versus batch size for the four primary implementations (end-to-end). Error bars are 95\,\% confidence intervals (smaller than the markers).}{fig:exec}

\colfigure{speedup_spartan_k40_native.pdf}{End-to-end speedup over the serial CPU ($T_{\mathrm{serial}}/T$, same batch size) of OpenMP (best thread count), the CUDA implementation with single-threaded preprocessing and the pipelined hybrid. The serial CPU is the baseline at 1.0.}{fig:speedup}

\colfigure{throughput_spartan_k40_native.pdf}{Throughput (images per second, log scale) of the four implementations versus batch size, end-to-end ($10{,}000/T$).}{fig:tput}

\subsection{CUDA Kernel Versus End-to-End Performance}\label{sec:compute}

This is the \emph{compute-only} comparison. Table~\ref{tab:compute} compares only the computation of the hidden layer, the output layer and the argmax; preprocessing and copies are excluded, so the GPU is represented by its kernels alone. For the CPU it is the sum of the timed stage durations. The GPU kernels are \val{Sp:comp:gpuvsserial:4096}$\times$ faster than serial at batch 4096 (\val{Sp:comp:kernel:4096}\,ms against \val{Sp:comp:serial:4096}\,ms) and \val{Sp:comp:gpuvsomp:4096}$\times$ faster than the best OpenMP configuration (\val{Sp:comp:omp:4096}\,ms). They beat OpenMP from batch 32 on (\val{Sp:comp:gpuvsomp:32}$\times$ at batch 32). For one image per batch the kernels are only \val{Sp:comp:gpuvsserial:1}$\times$ faster than serial and \val{Sp:comp:gpuvsomp:1}$\times$ as fast as OpenMP, i.e.\ slower: the batch gives a GPU too little parallel work.

Figure~\ref{fig:kernel} contrasts the two views for the CUDA implementation. The kernel-only speedup over the serial stages 1--3 is \val{Sp:comp:gpuvsserial:4096}$\times$ at batch 4096, about twice the \emph{end-to-end} speedup of \val{Sp:cuda:speedup:4096}$\times$, because the end-to-end time also contains the single-threaded preprocessing, the copies and the launch overhead (Section~\ref{sec:overhead}). The two ratios answer different questions and must not be mixed: the first compares only the neural-network computation, the second the complete pipeline an application waits for.

\begin{widetable}{Compute-only comparison on the Spartan node (stages 1--3, no preprocessing, no copies): time per pass (ms). OpenMP shows the best thread count in brackets; ``GPU kernels'' is the CUDA-event time of the three kernels (16$\times$16 block).}{tab:compute}
\begin{tabular}{rrrrrr}
\toprule
Batch & Serial & OpenMP (threads) & GPU kernels & \multicolumn{2}{c}{GPU kernels' speedup ($\times$)} \\
\cmidrule(lr){5-6}
 & (ms) & (ms) & (ms) & vs serial & vs OpenMP \\
\midrule
1 & 1{,}213.7 & 126.4 (16) & 1{,}091.0 & 1.11 & 0.12 \\
8 & 1{,}215.4 & 70.5 (24) & 143.9 & 8.45 & 0.49 \\
32 & 1{,}214.8 & 60.0 (24) & 47.3 & 25.67 & 1.27 \\
64 & 1{,}214.6 & 57.9 (24) & 34.4 & 35.28 & 1.68 \\
128 & 1{,}214.4 & 57.0 (24) & 29.8 & 40.75 & 1.91 \\
256 & 1{,}214.9 & 57.8 (24) & 28.1 & 43.29 & 2.06 \\
1024 & 1{,}214.8 & 56.7 (24) & 23.7 & 51.17 & 2.39 \\
4096 & 1{,}215.2 & 56.1 (24) & 23.0 & 52.77 & 2.44 \\
\bottomrule
\end{tabular}

\end{widetable}

\widefigure{cuda_kernel_vs_e2e_spartan_k40_native.pdf}{\cudabase{}: kernel-only time against end-to-end time (left, ms) and the corresponding speedups (right, $\times$): kernel-only against the serial time of stages 1--3 (compute-only), end-to-end against the serial end-to-end time. Kernel-only numbers exclude preprocessing, copies and launch overhead.}{fig:kernel}

\subsection{CUDA Transfer and Preprocessing Overhead (RQ4)}\label{sec:overhead}

Table~\ref{tab:breakdown} and Figure~\ref{fig:breakdown} split the pass of the CUDA implementation into stages. The two copies take \val{Sp:cuda:tshare:1}\,\% of the end-to-end time at batch 1 and \val{Sp:cuda:tshare:4096}\,\% at batch 4096. The host-to-device copy of the 31.4\,MB takes \val{Sp:cuda:h2d:4096}\,ms at batch 4096, which is \val{Sp:cuda:h2dgbs:4096}\,GB/s and agrees with the \val{Sp:micro:h2d_pageable_bandwidth:0}\,GB/s measured separately for pageable memory (Table~\ref{tab:micro}); pinned memory reaches \val{Sp:micro:h2d_pinned_bandwidth:0}\,GB/s.

At large batches the bigger non-GPU cost is the single-threaded preprocessing, which takes a constant \val{Sp:cuda:pre:4096}\,ms per pass. Its share of the end-to-end time grows from \val{Sp:cuda:preshare:1}\,\% at batch 1 to \val{Sp:cuda:preshare:4096}\,\% at batch 4096, while the kernels fall from \val{Sp:cuda:kshare:1}\,\% to \val{Sp:cuda:kshare:4096}\,\%. At batch 1 the launch, API and synchronization overhead (column ``Other'') is \val{Sp:cuda:otherms:1}\,ms per pass, \val{Sp:cuda:othershare:1}\,\% of the total, and it vanishes for large batches. This is why the end-to-end speedup stays well below the kernel-only speedup.

\colfigure{transfer_breakdown_spartan_k40_native.pdf}{Share of the end-to-end time of the CUDA implementation (single-threaded preprocessing) spent in each stage (\%) versus batch size.}{fig:breakdown}

\begin{widetable}{Time per pass of the CUDA implementation by stage on the Spartan node (ms). ``Other'' is the part of the wall-clock time not covered by the four timed stages (launch, API calls, synchronization).}{tab:breakdown}
\begin{tabular}{rrrrrrrr}
\toprule
Batch & Preprocess & H2D copy & Kernels & D2H copy & Other & Total & Copies share \\
 & (ms) & (ms) & (ms) & (ms) & (ms) & (ms) & (\% of total) \\
\midrule
1 & 20.4 & 85.4 & 1{,}091.0 & 213.3 & 96.8 & 1{,}507.0 & 19.8 \\
8 & 20.2 & 17.2 & 143.9 & 25.9 & 11.6 & 218.8 & 19.7 \\
32 & 20.2 & 7.8 & 47.3 & 7.0 & 3.1 & 85.5 & 17.3 \\
64 & 20.2 & 6.6 & 34.4 & 3.1 & 1.6 & 65.9 & 14.7 \\
128 & 20.2 & 6.0 & 29.8 & 1.8 & 0.8 & 58.6 & 13.2 \\
256 & 20.2 & 5.5 & 28.1 & 0.9 & 0.4 & 55.1 & 11.7 \\
1024 & 20.3 & 3.8 & 23.7 & 0.3 & 0.1 & 48.3 & 8.6 \\
4096 & 20.2 & 3.3 & 23.0 & 0.1 & 0.0 & 46.7 & 7.3 \\
\bottomrule
\end{tabular}

\end{widetable}

\begin{coltable}{Platform limits measured by micro-benchmarks.}{tab:micro}
\begin{tabular}{@{}lr@{}}
\toprule
Quantity & Measured value \\
\midrule
PCIe host-to-device, pageable / pinned & 9.5 / 10.5 GB/s \\
PCIe device-to-host, pageable / pinned & 10.0 / 10.5 GB/s \\
GPU device-to-device copy & 189.5 GB/s \\
Kernel launch, asynchronous / with sync & 3.12 / 7.44 $\mu$s \\
Xeon memory bandwidth, 1 / 8 / 24 threads & 11.8 / 46.5 / 43.4 GB/s \\
Xeon OpenMP fork--join, 8 / 24 threads & 2.13 / 7.16 $\mu$s \\
M4 memory bandwidth, 1 / 10 threads & 90.0 / 86.5 GB/s \\
M4 OpenMP fork--join, 4 / 10 threads & 7.25 / 21.98 $\mu$s \\
\bottomrule
\end{tabular}

\end{coltable}

\subsection{Hybrid OpenMP--CUDA Performance (RQ5)}\label{sec:hybrid}

Figure~\ref{fig:hybrid} and Table~\ref{tab:hybrid} compare OpenMP, the CUDA implementation, the hybrid-sequential control and the pipelined hybrid. The pipelined hybrid is faster than \cudabase{} at every batch size: \val{Sp:hyb:vscuda:1}$\times$ at batch 1, \val{Sp:hyb:vscuda:32}$\times$ at batch 32 and \val{Sp:hyb:vscuda:4096}$\times$ at batch 4096. Welch's $t$-test gives $p\le$ \val{Sp:test:hyb-cuda:maxp} for every batch size, and the 95\,\% bootstrap interval of the time ratio at batch 4096 is \val{Sp:test:hyb-cuda:cilo:4096} to \val{Sp:test:hyb-cuda:cihi:4096}. Against the best OpenMP configuration the hybrid is faster from batch \val{Sp:hyb:crossover} on, \val{Sp:hyb:vscpu:4096}$\times$ at batch 4096. At batch 1 and batch 8 OpenMP alone is faster (\val{Sp:omp:best:e2e:1}\,ms against \val{Sp:hyb:best:e2e:1}\,ms at batch 1). The CUDA pipeline with single-threaded preprocessing was not the fastest implementation at any batch size (Table~\ref{tab:best}).

Two ingredients contribute to the hybrid gain. Parallel preprocessing alone (the hybrid-sequential control) gives \val{Sp:hseq:vscuda:4096}$\times$ over \cudabase{} at batch 4096, and overlapping the batches adds another \val{Sp:hyb:vshseq:4096}$\times$. Even with a single preprocessing thread, pipelining with pinned memory gives \val{Sp:hyb:vscuda1:4096}$\times$. At batch 4096 the hybrid pass takes \val{Sp:hyb:best:e2e:4096}\,ms, only \val{Sp:hyb:vskernel:4096}$\times$ the kernel-only time of the CUDA implementation (\val{Sp:cuda:kernel:4096}\,ms): preprocessing and copies are almost completely hidden behind the kernels of other batches, so the GPU computation is again the bottleneck. Because the GPU stages overlap in the pipelined hybrid, only its end-to-end time is meaningful.

\widefigure{hybrid_compare_spartan_k40_native.pdf}{Left: end-to-end time per pass (ms, log scale) of OpenMP, the CUDA implementation with single-threaded preprocessing, the hybrid-sequential control and the pipelined hybrid (best thread counts). Right: speedup of the control and of the pipelined hybrid over the CUDA implementation ($\times$).}{fig:hybrid}

\begin{widetable}{Hybrid ablation on the Spartan node: time per pass (ms), speedup ($\times$) and significance. The first column is the CUDA implementation with one preprocessing thread; the hybrid-sequential column is a diagnostic control. Thread counts of the best configuration are in brackets.}{tab:hybrid}
\fitwide{\begin{tabular}{rrrrrrrr}
\toprule
Batch & CUDA-only & Hybrid-sequential (thr.) & Hybrid-pipelined (thr.) & Seq. vs CUDA-only & Pipelined vs CUDA-only & Pipelined vs OpenMP & Welch $p$ \\
 & (ms) & (ms) & (ms) & ($\times$) & ($\times$) & ($\times$) & (pipelined vs CUDA-only) \\
\midrule
1 & 1{,}507.0 & 1{,}520.6 (1) & 650.9 (1) & 0.99 & 2.32 & 0.25 & $8.7\times10^{-61}$ \\
8 & 218.8 & 207.5 (8) & 87.5 (8) & 1.05 & 2.50 & 0.86 & $2.4\times10^{-103}$ \\
32 & 85.5 & 69.9 (12) & 33.3 (2) & 1.22 & 2.57 & 1.86 & $3.9\times10^{-55}$ \\
64 & 65.9 & 49.9 (16) & 27.9 (24) & 1.32 & 2.36 & 2.12 & $1.6\times10^{-69}$ \\
128 & 58.6 & 41.9 (16) & 26.5 (24) & 1.40 & 2.21 & 2.20 & $1.2\times10^{-77}$ \\
256 & 55.1 & 38.2 (24) & 26.0 (24) & 1.44 & 2.12 & 2.27 & $5.9\times10^{-79}$ \\
1024 & 48.3 & 31.3 (12) & 25.2 (24) & 1.55 & 1.92 & 2.29 & $4.6\times10^{-33}$ \\
4096 & 46.7 & 28.8 (16) & 24.7 (24) & 1.62 & 1.89 & 2.32 & $4.8\times10^{-62}$ \\
\bottomrule
\end{tabular}
}
\end{widetable}

\begin{coltable}{Fastest implementation per batch size (end-to-end time per pass in ms; gap = second fastest / fastest). ``CUDA (1-thr. pre.)'' is the CUDA implementation with single-threaded preprocessing and a fixed 16$\times$16 block.}{tab:best}
\fitcol{\begin{tabular}{rlrlrr}
\toprule
Batch & Fastest & Time & Second fastest & Time & Gap \\
 & implementation & (ms) & implementation & (ms) & ($\times$) \\
\midrule
1 & OpenMP & 162.8 & Hybrid-pipelined & 650.9 & 4.00 \\
8 & OpenMP & 75.5 & Hybrid-pipelined & 87.5 & 1.16 \\
32 & Hybrid-pipelined & 33.3 & OpenMP & 61.9 & 1.86 \\
64 & Hybrid-pipelined & 27.9 & OpenMP & 59.2 & 2.12 \\
128 & Hybrid-pipelined & 26.5 & OpenMP & 58.2 & 2.20 \\
256 & Hybrid-pipelined & 26.0 & CUDA (1-thr. pre.) & 55.1 & 2.12 \\
1024 & Hybrid-pipelined & 25.2 & CUDA (1-thr. pre.) & 48.3 & 1.92 \\
4096 & Hybrid-pipelined & 24.7 & CUDA (1-thr. pre.) & 46.7 & 1.89 \\
\bottomrule
\end{tabular}
}
\end{coltable}

\subsection{CUDA Block Size and Tuning}\label{sec:block}

The block size of the hidden-layer kernel matters only for small batches (Table~\ref{tab:block}, Figure~\ref{fig:block}). At batch 64 the kernels take \val{Sp:cuda:blk8:kernel:64}\,ms with 8$\times$8 blocks, \val{Sp:cuda:blk16:kernel:64}\,ms with 16$\times$16 and \val{Sp:cuda:blk32:kernel:64}\,ms with 32$\times$32 blocks. At batch 4096 the three differ by less than 1\,\% (\val{Sp:cuda:blk8:kernel:4096}, \val{Sp:cuda:blk16:kernel:4096} and \val{Sp:cuda:blk32:kernel:4096}\,ms). A 32$\times$32 block at batch 32 launches only 8 thread blocks for the 15 multiprocessors of the K40m, consistent with its slower result (\val{Sp:cuda:blk32:kernel:32}\,ms against \val{Sp:cuda:blk8:kernel:32}\,ms for 8$\times$8). If the best of the three block sizes were used for the end-to-end time of the CUDA implementation, it would improve by at most \val{Sp:cuda:blockgain:max}\,\% (at batch \val{Sp:cuda:blockgain:argb}), and the batch size at which it catches up with the best OpenMP configuration would move from 256 to \val{Sp:cuda:bestblock:crossover} (where the two are within 0.1\,\%). The main comparison therefore keeps the fixed 16$\times$16 block and is labelled accordingly; it is not a best-versus-best comparison, because the hybrid was not block-tuned either.

\begin{widetable}{CUDA block-size comparison on the Spartan node: time per pass (ms) for the hidden-layer block sizes 8$\times$8, 16$\times$16 and 32$\times$32 (kernels only and end-to-end), the fastest block for the end-to-end time, and its gain over the fixed 16$\times$16 block. Block sizes of the output-layer and argmax kernels are fixed at 128 threads.}{tab:block}
\begin{tabular}{rrrrrrrcr}
\toprule
Batch & \multicolumn{3}{c}{CUDA kernels only (ms)} & \multicolumn{3}{c}{CUDA end-to-end (ms)} & Fastest & Gain over \\
\cmidrule(lr){2-4}\cmidrule(lr){5-7}
 & 8$\times$8 & 16$\times$16 & 32$\times$32 & 8$\times$8 & 16$\times$16 & 32$\times$32 & block (end-to-end) & 16$\times$16 (\%) \\
\midrule
1 & 1{,}095.6 & 1{,}091.0 & 1{,}136.9 & 1{,}510.1 & 1{,}507.0 & 1{,}560.6 & 16$\times$16 & 0.0 \\
8 & 142.7 & 143.9 & 144.8 & 218.3 & 218.8 & 219.6 & 8$\times$8 & 0.3 \\
32 & 45.3 & 47.3 & 56.7 & 83.2 & 85.5 & 94.8 & 8$\times$8 & 2.7 \\
64 & 37.3 & 34.4 & 48.7 & 69.1 & 65.9 & 81.1 & 16$\times$16 & 0.0 \\
128 & 29.5 & 29.8 & 34.6 & 58.2 & 58.6 & 63.3 & 8$\times$8 & 0.7 \\
256 & 26.5 & 28.1 & 29.6 & 53.6 & 55.1 & 56.6 & 8$\times$8 & 2.8 \\
1024 & 23.6 & 23.7 & 24.2 & 48.3 & 48.3 & 48.8 & 8$\times$8 & 0.0 \\
4096 & 23.1 & 23.0 & 23.2 & 46.8 & 46.7 & 46.8 & 16$\times$16 & 0.0 \\
\bottomrule
\end{tabular}

\end{widetable}

\colfigure{cuda_block_size_spartan_k40_native.pdf}{CUDA kernel time per pass (ms, log scale) for the thread-block sizes 8$\times$8, 16$\times$16 and 32$\times$32 of the hidden-layer kernel.}{fig:block}

\subsection{Initialization and Cold Start}\label{sec:cold}

The one-time costs that are excluded from the steady-state times were measured separately in 10 fresh processes (batch 256, 24 threads; Table~\ref{tab:cold} and Figure~\ref{fig:cold}). Creating the CUDA context takes \val{Sp:cold:init:mean}\,ms on average (median \val{Sp:cold:init:median}\,ms, range \val{Sp:cold:init:min}--\val{Sp:cold:init:max}\,ms), the weight upload \val{Sp:cold:upload:mean}\,ms and the buffer, stream and event allocation \val{Sp:cold:alloc:mean}\,ms, in total \val{Sp:cold:ctor:mean}\,ms, which is \val{Sp:cold:ctorvspass}$\times$ the time of one steady-state pass of the CUDA implementation. After this set-up the first CUDA pass takes \val{Sp:cold:cuda1:mean}\,ms and later passes \val{Sp:cold:cuda3:mean}\,ms, so the GPU kernels show no first-pass penalty. The first pipelined-hybrid pass is \val{Sp:cold:hyb:firstpenalty}\,\% slower than the second, and the first 24-thread OpenMP pass is \val{Sp:cold:omp:firstpenalty}\,\% slower (\val{Sp:cold:omp1:mean}\,ms against \val{Sp:cold:omp2:mean}\,ms); the serial code shows no difference. A program that processes a single pass of 10{,}000 images would therefore pay about \val{Sp:cold:ctor:mean}\,ms of GPU initialization on top of its pass; the steady-state numbers apply to programs that amortize this cost.

\widefigure{cold_start.pdf}{Cold start versus steady state on the Spartan node (10 fresh processes, batch 256, 24 OpenMP threads, mean with standard deviation). (a)~Time of the first and the second pass after the one-time GPU set-up. (b)~Components of the one-time GPU set-up, which is excluded from all steady-state times, against one steady-state CUDA pass.}{fig:cold}

\begin{widetable}{Cold-start measurements on the Spartan node (10 fresh processes, batch 256, 24 OpenMP threads): one-time GPU set-up and the first passes against later passes.}{tab:cold}
\begin{tabular}{lrrrrr}
\toprule
Quantity & Mean & Median & Std. dev. & Min & Max \\
 & (ms) & (ms) & (ms) & (ms) & (ms) \\
\midrule
CUDA context creation & 106.83 & 104.64 & 8.11 & 103.19 & 129.79 \\
Weight upload & 0.40 & 0.40 & 0.01 & 0.38 & 0.42 \\
Buffer, stream and event allocation & 2.97 & 2.94 & 0.08 & 2.88 & 3.11 \\
Total one-time GPU set-up & 110.65 & 108.53 & 8.07 & 106.93 & 133.49 \\
CUDA pass, 1st & 54.99 & 54.96 & 0.24 & 54.67 & 55.35 \\
CUDA pass, 3rd & 55.02 & 55.05 & 0.12 & 54.77 & 55.17 \\
Hybrid pass, 1st & 27.35 & 27.47 & 0.36 & 26.79 & 27.78 \\
Hybrid pass, 2nd & 25.95 & 25.93 & 0.12 & 25.80 & 26.19 \\
OpenMP (24 threads) pass, 1st & 80.84 & 80.45 & 1.55 & 79.90 & 85.20 \\
OpenMP (24 threads) pass, 2nd & 63.22 & 63.25 & 0.23 & 62.75 & 63.49 \\
Serial pass, 1st & 1{,}231.70 & 1{,}231.46 & 1.13 & 1{,}230.34 & 1{,}234.09 \\
Serial pass, 2nd & 1{,}232.32 & 1{,}232.13 & 0.96 & 1{,}231.11 & 1{,}234.06 \\
\bottomrule
\end{tabular}

\end{widetable}



\section{Discussion}
\subsection{Amdahl's law and parallel overhead (RQ6)}
We fitted $S(p)=1/(f+(1-f)/p)$ to the measured speedups of each batch size (Figure~\ref{fig:amdahl}). On the Xeon the fitted serial fraction is only $f=$ \val{Sp:amdahl:f:64} for batch 64 and \val{Sp:amdahl:f:4096} for batch 4096. This would allow at most \val{Sp:amdahl:ceil:64}$\times$ and \val{Sp:amdahl:ceil:4096}$\times$, so the measured \val{Sp:omp:speedup:4096:24}$\times$ at 24 threads is far from the Amdahl ceiling and the serial fraction is not what limits OpenMP here. The Karp--Flatt metric agrees: at batch 4096 it falls from \val{Sp:kf:4096:2} with 2 threads to \val{Sp:kf:4096:24} with 24 threads, so there is no growing overhead term. For one image per batch the fit gives $f=$ \val{Sp:amdahl:f:1}, a ceiling of only \val{Sp:amdahl:ceil:1}$\times$, and the Karp--Flatt metric stays between \val{Sp:kf:1:8} and \val{Sp:kf:1:12}: shrinking the batch increases the effective serial fraction, as Gustafson~\cite{gustafson1988reevaluating} describes for the reverse case of growing the problem. On the M4 the Karp--Flatt metric at batch 4096 jumps from \val{Mac:kf:4096:4} with 4 threads to \val{Mac:kf:4096:6} with 6 threads, where the efficiency cores join in. A heterogeneous chip violates the equal-processor assumption behind Amdahl's law, so the M4 value $f=$ \val{Mac:amdahl:f:4096} (R$^2=$ \val{Mac:amdahl:r2:4096}) is only a summary of the curve.

\widefigure{amdahl_combined.pdf}{Measured speedups (dots) and fitted Amdahl curves for four batch sizes on the Apple M4 (left) and the Xeon (right); $f$ is the fitted serial fraction.}{fig:amdahl}

Synchronisation explains why OpenMP loses for single images on the M4 and wins on the Xeon. One image needs \val{Sp:serial:usimg:1}\,$\mu$s of serial work on the Xeon but only \val{Mac:serial:usimg:1}\,$\mu$s on the M4, while one OpenMP fork--join costs \val{Sp:micro:openmp_fork_join:24}\,$\mu$s with 24 threads on the Xeon and \val{Mac:micro:openmp_fork_join:10}\,$\mu$s with 10 threads on the M4 (Table~\ref{tab:micro}). When the work per batch is comparable to the cost of synchronising the threads, as on the M4, parallelisation loses; when it is much larger, as on the Xeon, it wins even for one image. A larger batch spreads the synchronisation cost over more images, which is why OpenMP improves with the batch size.

\subsection{Memory traffic and the GPU}
If every image had to read the 803{,}840-byte hidden-layer weights from device memory, one pass would move \val{X:weights:GB}\,GB. At the \val{Sp:micro:d2d_copy_bandwidth:0}\,GB/s measured with a device-to-device copy this would take about 42\,ms, longer than the measured kernel time of \val{Sp:cuda:kernel:4096}\,ms at batch 4096; the weights must therefore be served mostly from on-chip caches, which is the data reuse that the roofline model~\cite{williams2009roofline} makes visible. The PCIe link is much slower than device memory (\val{Sp:micro:h2d_pinned_bandwidth:0}\,GB/s pinned against \val{Sp:micro:d2d_copy_bandwidth:0}\,GB/s), so data have to be reused on the device to pay for the transfer; with 3{,}136 bytes per image crossing the link once, the copy costs only \val{Sp:cuda:h2d:4096}\,ms per pass here~\cite{gregg2011data,harris2012transfers}. On the CPU side, the serial code also re-reads the weights for every image: this is \val{Sp:serial:wtraffic:4096}\,GB/s of weight traffic, below the \val{Sp:micro:cpu_triad_bandwidth:1}\,GB/s one thread achieves on the memory benchmark, and compiling with \texttt{-march=native} makes it only \val{X:serial:O3OverNative:4096}$\times$ faster than plain \texttt{-O3} (Table~\ref{tab:flags}). Both observations are consistent with a loop limited by data movement rather than arithmetic. With 24 threads the weight traffic is \val{Sp:omp:best:wtraffic:4096}\,GB/s, far above the memory bandwidth of \val{Sp:micro:cpu_triad_bandwidth:24}\,GB/s measured with 24 threads, so the weights are served from the caches there as well.

\subsection{Launch overhead and GPU utilisation}
An asynchronous kernel launch costs \val{Sp:micro:kernel_launch_async:0}\,$\mu$s and a launch followed by a synchronisation \val{Sp:micro:kernel_launch_sync:0}\,$\mu$s (Table~\ref{tab:micro}). For one image per batch the program spends \val{Sp:cuda:otherus:1}\,$\mu$s per batch on launches, API calls and synchronisation, but the kernels themselves take \val{Sp:cuda:kus:1}\,$\mu$s per batch. With this mapping a batch of one image has 256 useful threads out of 4{,}096 launched (16$\times$16 blocks), so most of the K40m's multiprocessors are idle and every thread executes a 784-step dependent chain. This is a latency-bound situation that a larger batch removes, and the kernel time per pass falls from \val{Sp:cuda:kernel:1}\,ms at batch 1 to \val{Sp:cuda:kernel:4096}\,ms at batch 4096.

\subsection{Which implementation is best}
Table~\ref{tab:best} summarises the measurements. For one image per batch and for batch 8, OpenMP on the multicore CPU is fastest. From batch 32 on, the pipelined hybrid is fastest, and its advantage over the second-best option is between \val{Sp:hyb:vscpu:32}$\times$ (batch 32) and about 2$\times$ at larger batches. \cudabase{} is never the fastest: its single-threaded preprocessing and its synchronous copies leave the GPU idle for part of every batch, which the hybrid removes. These rankings hold for both compiler settings (the \texttt{-O3} runs show the same order), but the batch size at which \cudabase{} (fixed 16$\times$16 block) overtakes the best OpenMP configuration moves from \val{Sp:cuda:crossover} (\texttt{-march=native}) to \val{SpO3:cuda:crossover} (plain \texttt{-O3}), because the preprocessing stage is faster with plain \texttt{-O3} (\val{SpO3:cuda:pre:4096}\,ms against \val{Sp:cuda:pre:4096}\,ms per pass) and the CPU baselines are slightly slower (Table~\ref{tab:flags}).

\begin{coltable}{Effect of the compiler setting on the Spartan node: time per pass (ms) with \texttt{-march=native} and with plain \texttt{-O3}, and their ratio.}{tab:flags}
\fitcol{\begin{tabular}{lrrrrrr}
\toprule
Implementation & \multicolumn{3}{c}{Batch 256} & \multicolumn{3}{c}{Batch 4096} \\
\cmidrule(lr){2-4}\cmidrule(lr){5-7}
 & native (ms) & \texttt{-O3} (ms) & ratio & native (ms) & \texttt{-O3} (ms) & ratio \\
\midrule
Serial & 1{,}235.2 & 1{,}296.2 & 1.05 & 1{,}235.4 & 1{,}295.6 & 1.05 \\
OpenMP, best thread count & 58.9 & 61.1 & 1.04 & 57.3 & 61.1 & 1.07 \\
CUDA-only (16$\times$16) & 55.1 & 51.4 & 0.93 & 46.7 & 42.7 & 0.91 \\
Hybrid-pipelined, best & 26.0 & 26.1 & 1.01 & 24.7 & 24.8 & 1.00 \\
\bottomrule
\end{tabular}
}
\end{coltable}


\section{Limitations and Threats to Validity}\label{sec:limits}
\textbf{Hardware and software.} One GPU and one CPU model were measured, on one node type. The Tesla K40m is a Kepler-generation GPU (compute capability 3.5, which \texttt{nvcc} 11.7 flags as deprecated), so it is not representative of recent GPUs. Newer GPUs with faster links and more memory bandwidth may move the crossover points; we did not test them. Only g++ 11.4.1, Apple clang 21.0.0 and CUDA 11.7 were used.

\textbf{Workload.} The network is small: its \val{X:weights:GB}\,GB of weight traffic per pass is served from caches, and a larger or convolutional network whose weights do not fit in cache may be limited by memory bandwidth, which would change the CPU--GPU balance. Only float32 inference was measured, with no energy or power measurements.

\textbf{Implementation quality.} The loops and kernels are straightforward: the CPU code is blocked only to allow vectorisation and the CUDA kernel uses coalesced loads without shared-memory tiling; no BLAS library is used. The effective rates at batch 4096 (\val{Sp:serial:gflops:4096}\,GFLOP/s serial, \val{Sp:omp:best:gflops:4096}\,GFLOP/s OpenMP, \val{Sp:cuda:kgflops:4096}\,GFLOP/s for the CUDA kernels alone) show that none of them is tuned. Tuned kernels would change the absolute speedups and possibly the crossover batch sizes; our conclusions concern straightforward implementations. We did not use hardware peak numbers because we did not measure them.

\textbf{Design choices that affect the comparison.} The preprocessing stage is part of our workload definition, and the hybrid gain depends on it; the CUDA implementation deliberately uses one preprocessing thread, which is why the hybrid-sequential control is reported. OpenMP and the pipelined hybrid use the best of seven thread counts per batch size, whereas CUDA (and the hybrid's kernels) use a fixed 16$\times$16 block; the comparison is therefore ``CUDA with a fixed block against OpenMP with tuned thread count'', not best-versus-best. The best block would improve the CUDA end-to-end time by at most \val{Sp:cuda:blockgain:max}\,\% (Section~\ref{sec:block}). One-time GPU initialisation, allocation and weight upload are excluded from the steady-state times, which favours the GPU for short-lived programs; they cost \val{Sp:cold:ctor:mean}\,ms in total (Section~\ref{sec:cold}).

\textbf{Measurement noise.} The Spartan node is managed by a shared scheduler (our jobs held 24 of its 28 CPUs and one of its two GPUs, so other jobs could use the rest), and the M4 ran with background services (system load average about 3 before the runs). Run-to-run variation is small: the median coefficient of variation is \val{Sp:cv:median}\,\% on the Spartan node and \val{Mac:cv:median}\,\% on the M4 (90th percentile \val{Sp:cv:p90}\,\% on Spartan). The largest variation on Spartan is the \val{Sp:slow:impl} configuration at batch \val{Sp:slow:batch} with \val{Sp:slow:threads} threads (\val{Sp:slow:n} passes): mean \val{Sp:slow:mean}\,ms, median \val{Sp:slow:median}\,ms, standard deviation \val{Sp:slow:std}\,ms, minimum \val{Sp:slow:min}\,ms and maximum \val{Sp:slow:max}\,ms, a coefficient of variation of \val{Sp:slow:cv}\,\%. Exactly \val{Sp:slow:count} pass exceeds 1.5 times the median; it is kept in all statistics and in the raw data. On the M4 the largest variation is the serial run at batch \val{Mac:slow:batch}: mean \val{Mac:slow:mean}\,ms, median \val{Mac:slow:median}\,ms, standard deviation \val{Mac:slow:std}\,ms, minimum \val{Mac:slow:min}\,ms and maximum \val{Mac:slow:max}\,ms (\val{Mac:slow:cv}\,\%). The 20 passes of a configuration run back to back, so they are not independent samples of the machine state. macOS ignores thread binding, the Xeon node has two NUMA nodes (threads above 14 cross a socket boundary), and the M4's two core types violate the assumptions of Amdahl's law. The confidence intervals therefore describe the repeatability of these runs, not all possible conditions.

\textbf{Pipelined hybrid.} CUDA events of concurrent streams include time spent waiting for other streams, so the stage durations of the pipelined hybrid are not true durations; we report only its end-to-end time.


\section{Conclusion and Future Work}
We implemented one MNIST inference workload four ways---serial, OpenMP, CUDA, and a hybrid OpenMP--CUDA pipeline---verified that they agree (largest logit difference \val{Sp:val:cuda:maxerr}, identical predictions), and measured them with \val{Sp:runs} passes per configuration. The main findings are:

\begin{itemize}
\item \textbf{RQ1.} OpenMP reaches \val{Sp:omp:speedup:4096:24}$\times$ over serial with 24 threads on the Xeon (batch 4096, efficiency \val{Sp:omp:eff:4096:24}), but only \val{Mac:omp:speedup:4096:10}$\times$ with 10 threads on the M4, and it loses to serial for single images on the M4.
\item \textbf{RQ2, RQ3.} CUDA is \val{Sp:cuda:speedup:4096}$\times$ faster than serial end-to-end at batch 4096 but \val{Sp:cuda:speedup:1}$\times$ (slower) at batch 1; its kernel-only speedup (\val{Sp:cuda:kspeedup:4096}$\times$) is twice the end-to-end speedup. It overtakes the best OpenMP configuration only for batches of about 128 to 256 images or more.
\item \textbf{RQ4.} The copies take \val{Sp:cuda:tshare:4096}\,\% to \val{Sp:cuda:tshare:1}\,\% of the CUDA time; the single-threaded preprocessing (up to \val{Sp:cuda:preshare:4096}\,\%) is the larger overhead for big batches.
\item \textbf{RQ5.} The pipelined hybrid needs \val{Sp:hyb:best:e2e:4096}\,ms at batch 4096, \val{Sp:hyb:vscuda:4096}$\times$ faster than CUDA-only and \val{Sp:hyb:vscpu:4096}$\times$ faster than the best OpenMP; it is fastest from batch \val{Sp:hyb:crossover} on, while OpenMP is fastest for batches of 1 and 8.
\item \textbf{RQ6.} The fitted Amdahl serial fraction is \val{Sp:amdahl:f:4096} at batch 4096 on the Xeon but \val{Sp:amdahl:f:1} for single images, and synchronisation cost against work per batch explains the different behaviour of the Xeon and the M4.
\end{itemize}

The answer to ``which is best'' therefore depends on the batch size: multicore OpenMP for tiny batches, and a hybrid that overlaps CPU preprocessing with GPU transfer and compute for larger ones; a GPU used naively (synchronous, single-threaded preprocessing) is never the best choice here. Future work includes tiled and library-based kernels (cuBLAS), newer GPUs, larger and convolutional networks, GPU-side preprocessing, splitting one batch across CPU and GPU, and energy measurements.

\textbf{Availability.} The code, raw measurements, analysis scripts and the sources of this paper are in a private repository and in the accompanying Overleaf project (folder \texttt{code}).


%% Text printed under the heading of the reference section (see \csReferencesIntro in cs259.cls).
\renewcommand{\csReferencesIntro}{%
  \noindent\emph{Original material.} All figures and tables in this paper were created by the author for this paper from the
  project's own measurements, source code, and calculations; none is copied or adapted from another publication. The data figures
  are generated by scripts from the raw measurements, and the diagrams (workload flow chart, execution architecture, and hybrid time line)
  were drawn by the author. The source code (serial, OpenMP, CUDA, and hybrid implementations, benchmark harness, and analysis scripts),
  the raw measurements, and the report source are the author's own work and are listed as reference~\cite{induru2026repository}.
  All other entries are published works cited in the text. The references follow the IEEE style; as the course template requires,
  they are numbered and ordered alphabetically by first author.\par\smallskip}

\bibliographystyle{cs259ieee}
\bibliography{refs}

\end{document}
